\documentclass[a4paper,oneside, 11pt,english,listof=totoc,bibliography=totoc]{scrbook}

% recommended defaults
\usepackage{babel}
\usepackage[utf8]{inputenc}
\usepackage[T1]{fontenc}
\usepackage{dsfont}
\usepackage{siunitx}
\usepackage{multirow}
\sisetup{locale=DE,per-mode=fraction,separate-uncertainty=true}
% references, use biber
\usepackage[style=numeric,sorting=none,backend=biber]{biblatex}
\addbibresource{refs.bib}
% optional packages, yet recommended
\usepackage{amsmath,amssymb,amsthm}
\usepackage{mathdots}
\usepackage{subcaption}
\usepackage{hyperref}
\usepackage{cleveref}
\usepackage{microtype}
\usepackage{booktabs}
\usepackage{csquotes}
\usepackage{placeins}
\usepackage{marvosym}
\usepackage{tabularx}
\usepackage{rotating}

\usepackage{graphicx}
\graphicspath{%
  {figs/}%
}
\usepackage{pdfpages}
\usepackage{xspace}
\usepackage{xcolor}
\usepackage{mdframed}

\title{Determining the Axial L{\"u}scher Term of an $SU(2)$ Quantum Link Model}
\author{Josh Lengel}
\date{}

\begin{document}

\frontmatter{}

\maketitle

\tableofcontents

\mainmatter{}

\chapter{Introduction}
This paper aims to investigate the axial L{\"u}scher term of a two-dimensional SU(2) quantum link model. Using numerical methods such as tensor networks and density matrix renormalization group (DMRG) on a hexagonal lattice, the quark-antiquark-potential will be computed for different coupling strengths $g^2$. Finally, the shape of this potential will be analyzed to determine the existence and behavior of the L\"uscher term.

\section{Wilson Theory}
Beginning with Wilson's formulation of lattice gauge theory\cite{Gattringer}, the two-dimensional space is discretized into a lattice with sites labeled by $x$ and link variables $u_{x,i}$ connecting site $x$ to one of its neighbors $x+\hat{i}$ in the direction labeled by $i$. Link variables are related to the gauge field $A^\mu$ of the standard Yang-Mills theory by $u_{x,k}=e^{iaA^k\left(x+\hat{k}/2\right)}$.

The link variables $u_{x,i}$ are elements of the gauge group instead of the corresponding Lie algebra and transform as $u'_{x,i}=V_{x}u_{x,i}V^{-1}_{x+\hat{i}}$ under local gauge transformations $V_x$. 

\section{Hamiltonian Formulation}
Consider the Kogut-Susskind Hamiltonian formulation of Wilson's lattice gauge theory,
\begin{equation}\label{eq:hamiltonian}
    \hat{H}=g^2\sum\limits_{x,i}\left({\hat{\vec{L}}}^2_{x,i} + {\hat{\vec{R}}}^2_{x,i}\right)+\frac{1}{g^2}\sum\mathrm{Tr}\left(U_\mathrm{p}+U^\dagger_\mathrm{p}\right)
\end{equation}
whereby $g^2$ is the squared coupling constant, ${\hat{\vec{L}}}^2_{x,i}$ and ${\hat{\vec{R}}}^2_{x,i}$ are the electric (momentum) operators and $U_\mathrm{p}$ is defined as the product of link variables in a counter-clockwise direction around the smallest possible loop (plaquette). Similarly, $U^\dagger_\mathrm{p}$ represents the product of the inverse of link variables around a plaquette in the clockwise direction\cite{KogutSusskind}.

The electric operators are defined by the commutation relations
\begin{equation}\label{eq:lr-commutation}
    \begin{array}{rl}
        \left[L^a_{x,i},U_{y,j}\right]&=-\delta_{xy}\delta{ij}\frac{\tau^a}{2}U_{y,j}, \\
        \left[R^a_{x,i},U_{y,j}\right]&=\delta_{xy}\delta{ij}U_{y,j}\frac{\tau^a}{2}, \\
        \left[L^a_{x,i},L^b_{y,j}\right]&=i\delta_{xy}\delta_{ij}\epsilon_{abc}L^c_{x,i}, \\
        \left[R^a_{x,i},R^b_{y,j}\right]&=i\delta_{xy}\delta_{ij}\epsilon_{abc}R^c_{x,i}, \\
        \left[L^a_{x,i},R^b_{y,j}\right]&=0
    \end{array}
\end{equation}
whereby $U_{x,i}$ represents the operator corresponding to the link variable $u_{x,i}$ and $\tau^a$ represents the Pauli matrices acting on color space.

Gauge invariance requires that Gauss' law
\begin{equation}\label{eq:gauss}
    G^a_x\left|\psi\right\rangle=\sum\limits_{i}\left(L^a_{x,i}+R^a_{x-\hat{i},i}\right)\left|\psi\right\rangle=0
\end{equation}
hold for any element $\psi$ of the Hilbert space of wavefunctions. Since $G^a_x$ is a generator of infinitesimal gauge transformations, it must commute with the Hamiltonian and be constant. More specifically, the spectrum of $G^a_x$ represents conserved charges, of which there are none in this model.

Now, the primary issue with this approach is that since $[U_{x,i}, U_{y,j}]=0$, the Hilbert space of a lattice of any size is infinite-dimensional. This means that the Hamiltonian must be truncated in order to numerically diagonalize it. However, truncating the representation typically affects the underlying gauge symmetry if done poorly.

\section{Quantum Link Model}
At this point, introducing a quantum link model will provide a method for truncation which preserves the underlying gauge symmetry of the theory. The SU(2) link variables $u_{i,x}$ and corresponding operators will be substituted with the quantum link operators
\begin{equation}\label{eq:qlo}
    U_{i,x}=U^0_{i,x}\mathds{1}+U^a_{i,x}\tau^a
\end{equation}
with the following, additional commutation relations
\begin{equation}\label{eq:u-commutation}
    \begin{array}{rl}
        \left[U^0_{x,i},U^0_{y,j}\right]&=0, \\
        \left[U^0_{x,i},U^a_{y,j}\right]&=i\delta_{xy}\delta{ij}\left(R^a_{x,i}-L^a_{x,i}\right), \\
        \left[U^a_{x,i},U^b_{y,j}\right]&=i\delta_{xy}\delta{ij}\epsilon_{abc}\left(R^c_{x,i}+L^c_{x,i}\right)
    \end{array}
\end{equation}
thereby embedding the $\mathfrak{su}(2)\times\mathfrak{su}(2)$ link-algebra in a larger $\mathfrak{so}(5)$ algebra. The $R^a_{x,i}$ and $L^a_{x,i}$ operators consist of three components each and, including the four operators $U^\rho_{x,i}$, represent the ten generators of $SO(5)$. By construction, this quantum link model still retains the gauge invariance and the Wilson theory may be recovered by dimensional reduction\cite{Chandrasekharan, Beard}.

$\mathfrak{so}(5)$ has many finite representations to choose from, the simplest of which that is non-trivial being the four-dimensional spinor representation. This representation is not ideal however, since link states are not symmetric and feature different $SU(2)$ elements on either end. The five-dimensional vector representation will be used from here on out since the link states are symmetric for each of the five basis states.

The process of promoting link variables to non-commuting operators is motivated by similar approaches to modelling spin interaction. In it's classical formulation, the Heisenberg spin model,
\begin{equation*}
    H=-\sum\limits_{i=1}^{N-1}\vec{s}_i\cdot\vec{s}_{i+1}
\end{equation*}
depends on the spin vectors of each site $i$. Since $\vec{s}_i$ can take on infinitely many values, a probability distribution over the configuration space would be infinite-dimensional. This motivates promoting the spin vectors $\vec{s}_i$ to spin operators $\hat{\vec{S}}_i$ with the commutation relation
$$[\hat{S}_{x,i},\hat{S}_{y,j}]=\delta_{xy}\epsilon_{ijk}S_{x,k}.$$
Importantly, the $O(3)$ symmetry is still preserved by construction. The resulting operator algebra is now a product basis of the two-dimensional $\mathfrak{su}(2)$ which results in a finite-dimensional Hilbert space.

\section{Rishon Representation}
Next, each link will be represented by two fermions, known as rishons, on either end\footnote{Note that these have no physical significance whatsoever, they are merely a mathematical tool to construct a simpler basis for the Hilbert space.}. These are represented by the pair of fermionic creation and annihilation operators
\begin{equation}\label{eq:rishon}
    \left(c^{a\dagger}_{x,i,S}, c^a_{x,i,S}\right)
\end{equation}
for which the stardard anticommutation relations
\begin{equation}\label{eq:ac-rishon}
    \begin{array}{l}
        \left\{c^{a\dagger}_{x,i,S_1},c^b_{y,j,S_2}\right\}=\delta^{ab}\delta_{xy}\delta_{ij}\delta_{S_1S_2} \\
        \left\{c^a_{x,i,S_1},c^b_{y,j,S_2}\right\}=0 \\
        \left\{c^{a\dagger}_{x,i,S_1},c^{b\dagger}_{y,j,S_2}\right\}=0
    \end{array}
\end{equation}
hold. This ensures that either zero or one of each type of rishon may exist at each end of the link.

These rishons come in two types, which will be labeled by $a\in\{1,2\}$ and are placed at the end $S\in\{+,-\}$ of the link $x,i$. All of the $SO(5)$ generators can now be restated in terms of these operators:
\begin{align*}
    L^1_{x,i}&=\frac{c^{1\dagger}_{x,i,-}c^2_{x,i,-}+c^{2\dagger}_{x,i,-}c^1_{x,i,-}}{2} \\
    L^2_{x,i}&=i\frac{-c^{1\dagger}_{x,i,-}c^2_{x,i,-}+c^{2\dagger}_{x,i,-}c^1_{x,i,-}}{2} \\
    L^3_{x,i}&=\frac{c^{1\dagger}_{x,i,-}c^1_{x,i,-}-c^{2\dagger}_{x,i,-}c^2_{x,i,-}}{2} \\
    R^1_{x,i}&=\frac{c^{1\dagger}_{x,i,+}c^2_{x,i,+}+c^{2\dagger}_{x,i,+}c^1_{x,i,+}}{2} \\
    R^2_{x,i}&=i\frac{-c^{1\dagger}_{x,i,+}c^2_{x,i,+}+c^{2\dagger}_{x,i,+}c^1_{x,i,+}}{2} \\
    R^3_{x,i}&=\frac{c^{1\dagger}_{x,i,+}c^1_{x,i,+}-c^{2\dagger}_{x,i,+}c^2_{x,i,+}}{2}
\end{align*}

The expressions for $L^3_{x,i}$ and $R^3_{x,i}$ allow the rishons of type $a=1,2$ to be identified with color charges of $+\frac{1}{2}$ ($\uparrow$) and $-\frac{1}{2}$ ($\downarrow$), respectively. For the quantum link operators one obtains the representation
\begin{equation*}
    U_{x,i}=\begin{pmatrix}
        c^{2\dagger}_{x,i,-}c^2_{x,i,+}+c^{1\dagger}_{x,i,+}c^1_{x,i,-} & -c^{2\dagger}_{x,i,-}c^1_{x,i,+}+c^{2\dagger}_{x,i,+}c^1_{x,i,-} \\
        -c^{1\dagger}_{x,i,-}c^2_{x,i,+}+c^{1\dagger}_{x,i,+}c^2_{x,i,-} & c^{1\dagger}_{x,i,-}c^1_{x,i,+}+c^{2\dagger}_{x,i,+}c^2_{x,i,-}
    \end{pmatrix}
\end{equation*}
Since each term contains the annihilation of a rishon on one end and the creation of a rishon on the other, these operators have the effect of transferring charges across the link. Furthermore, each of the ten operators conserves the number of rishons on each link. It can be shown that the number of rishons on each link corresponds directly to the choice of representation of $SO(5)$. For the five-dimensional representation, there are always two rishons on each link. This leads to five basis states, four of which are simply the combinations of one up-type or down-type rishon at each end. The fifth is more subtle but represents the superposition of states where two rishons of opposite type are created at the same end of the link, forming a singlet.
\begin{align*}
    |1\rangle&=|\downarrow\rangle|\downarrow\rangle \\
    |2\rangle&=|\uparrow\rangle|\downarrow\rangle \\
    |3\rangle&=|\downarrow\rangle|\uparrow\rangle \\
    |4\rangle&=|\uparrow\rangle|\uparrow\rangle \\
    |5\rangle&=\frac{1}{2}\left(|\uparrow\downarrow\rangle|\;\rangle+\rangle|\;\rangle|\uparrow\downarrow\right)
\end{align*}

\section{Even Site Basis}\label{sec:site-basis}
In order to simplify this representation further, the gauge-specific information contained by this representation will be removed. Gauss' law (see \ref{eq:gauss}) ensures that no state is valid where an odd number of rishons of any type meet at a lattice site.  Furthermore, applying Gauss' law leads to only four valid states at any given lattice site.
\begin{align*}
    &1: & &|\uparrow\downarrow\rangle_1\times|\uparrow\downarrow\rangle_2\times|\uparrow\downarrow\rangle_3 \\
    &2: & &\frac{1}{2}\left(|\uparrow\downarrow\rangle_1\times|\uparrow\rangle_2\times|\downarrow\rangle_3-|\uparrow\downarrow\rangle_1\times|\downarrow\rangle_2\times|\uparrow\rangle_3\right)\\
    &3: & &\frac{1}{2}\left(|\uparrow\downarrow\rangle_2\times|\uparrow\rangle_3\times|\downarrow\rangle_1-|\uparrow\downarrow\rangle_2\times|\downarrow\rangle_3\times|\uparrow\rangle_1\right)\\
    &4: & &\frac{1}{2}\left(|\uparrow\downarrow\rangle_3\times|\uparrow\rangle_1\times|\downarrow\rangle_2-|\uparrow\downarrow\rangle_3\times|\downarrow\rangle_1\times|\uparrow\rangle_2\right)
\end{align*}
whereby states 2/3/4 are permutations of each other. The $|\uparrow\downarrow\rangle$ state represents the singlet state with zero color charge. It should be noted here that these rishon states do not exist in isolation but are only one half of a link state. However, if the rishon states of three surrounding lattice sites is fixed, this construction uniquely determines the rishon state of the remaining site, up to a gauge transformation flipping up and down states. There is a very natural and simple way to express these states and simultaneously encode Gauss' law. By removing the type information of a rishon, the singlet state will now be represented as an empty state and any other state, whether up or down, will be represented as a dot $\bullet$. This is valid since the two rishons must contain opposite charges at the lattice site to fulfill Gauss' law and the choice of which rishon is up-type or down-type only results in a sign-flip, which does not change the probability amplitude or behavior of the state. The four allowed states are displayed with this representation in figure \ref{fig:site-basis}.
\begin{figure}
    \centering
    \includegraphics[width=0.7\linewidth]{figs/site-basis.png}
    \caption{Four basis states in the rishon representation. The gray arrow represents the sign of the singlet\cite{PLudwig}}
    \label{fig:site-basis}
\end{figure}

Finally the even site basis can be constructed. Using the four site basis states defined above, every even lattice site is represented as a local product basis of the four possible states. Gauss' law must still be applied at the odd sites however, although this only limits the number of possible states further.

\section{Plaquette Environments}
Now that a suitable basis has been chosen to represent the lattice, the magnetic Hamiltonian (see \ref{eq:hamiltonian}) can be reformulated in terms of this basis. For any given plaquette, the action of the magnetic Hamiltonian clearly preserves the external links connecting lattice sites that are part of the plaquette with those which are not. As such, the $4^3=64$ unique plaquette configurations can be assigned to eight distinct environments which are unique up to reflection and rotation. Each environment has two subtypes, which corresponds to replacing each rishon in the plaquette with a singlet and vice versa (see figure \ref{fig:environments}).
\begin{figure}
    \centering
    \includegraphics[width=0.9\linewidth]{figs/environments.png}
    \caption{All eight plaquette environments with both subtypes\cite{PLudwig}}
    \label{fig:environments}
\end{figure}
\clearpage
Since the Hamiltonian must respect these symmetries, the general form of the plaquette hamiltonian can be written as
\begin{equation}\label{eq:plaquette-hamiltonian}
    H_\mathrm{p}=\begin{pmatrix}
        H_\mathrm{I} & & & &\hdots & & & 0 \\
        & H_\mathrm{II} & & & & & & \\
        & & H_\mathrm{III} & & &\iddots & & \\
        & & & H_\mathrm{IV} & & & & \vdots \\
        \vdots& & & & H_\mathrm{V} & & & \\
        & &\iddots & & & H_\mathrm{VI} & & \\
        & & & & & & H_\mathrm{VII} & \\
        0& & &\hdots & & & & H_\mathrm{VIII}
    \end{pmatrix}
\end{equation}
whereby each of the $H_E$ is a block diagonal matrix with a 2x2 matrix $H_\mathrm{env}$ in each block acting on the two subtypes of each configuration belonging to environment $E$. Since $H$ is hermitian,
\begin{equation}\label{eq:e-hamiltonian}
    H_\mathrm{env}=\begin{pmatrix}
        W^1_\mathrm{env} & T_\mathrm{env} \\
        T^*_\mathrm{env} & W^2_\mathrm{env}
    \end{pmatrix}
\end{equation}

When the two subtypes are related by reflection or rotation, the symmetry argument demands
$$W^1_\mathrm{env}=W^2_\mathrm{env}\text{ and }T_\mathrm{env}=T^*_\mathrm{env}$$
Multiplying out the plaquette operator in the rishon representation and taking its non-zero matrix elements yields the following values for $W$ and $T$:
\begin{equation}\label{eq:WT-plaquette}
    W=0,\;T_\mathrm{I}=-128,T_\mathrm{II}=T_\mathrm{III}=T_\mathrm{IV}=32,T_\mathrm{V}=T_\mathrm{VI}=T_\mathrm{VII}=-32,T_\mathrm{VIII}=32
\end{equation}

Finally the magnetic Hamiltonian is given by the sum over all plaquettes
\begin{equation}\label{eq:mag-hamiltonian}
    H_\mathrm{mag}=\sum_p H_p
\end{equation}
whereby each term in the sum is really an outer product between identity elements from the left for all previous lattice sites and from the right for all the following lattice sites.

\section{Full Hamiltonian}
It is easy to see that the eigenvalue of each ${L^a_{x,i}}^2$ and ${R^a_{x,i}}^2$ is $1/4$ for a particular rishon in an up or down state. For the singlet state the eigenvalue is zero. It follows that each rishon contributes $6/4$ to the electric hamiltonian. Since rishons represent the entire link in the even site basis, this value must be doubled to account for the rishon at the other end of the link. Given a site basis of the four states $(\_,\_,\_)$, $(\bullet,\bullet,\_)$, $(\_,\bullet,\bullet)$ and $(\bullet,\_,\bullet)$, the electric Hamiltonian of an individual site is represented by the diagonal matrix
\begin{equation}\label{eq:el-hamiltonian}
    H_\mathrm{x,el}=\begin{pmatrix}
        0 & 0 & 0 & 0 \\
        0 & 2\cdot 6/4 & 0 & 0 \\
        0 & 0 & 2\cdot 6/4 & 0 \\
        0 & 0 & 0 & 2\cdot 6/4
    \end{pmatrix}
\end{equation}
The full electric Hamiltonian is recovered by combining the individual site Hamiltonians similarly to the magnetic Hamiltonian.

\section{Cornell Potential}
The allowed energies of the system are obtained via the eigenvalues of the Hamiltonian. By modifying Gauss' law at the top and bottom of the lattice (see \ref{sec:potential}) one can introduce two opposing color charges. The potential energy of a quark-antiquark pair held at a distance $r$ should follow
\begin{equation}\label{eq:cornell}
    V(r)=\frac{\gamma}{r}+\sigma r+c
\end{equation}
whereby $\gamma$ is the L{\"u}scher term, $\sigma$ is the string tension and $c$ is an arbitrary constant used for fitting this function to the data\cite{Cornell}. It becomes clear from this function that the L\"uscher term will only be observable for small $r$. This will become problematic as finite-size effects are large in this domain.

\chapter{Methods}
Now that the Hamiltonian has been constructed in an appropriate basis, the solution involves determining its eigenvectors and -values. However, given a lattice with $N$ sites, the Hilbert space is a product space of $N/2$ sites, meaning it is $4^{N/2}$-dimensional! This is much too large for direct diagonalization and the asymptotic runtime would be dreadful.

\section{Tensor Networks}
To solve this problem, the full $4^{N/2}$-dimensional vector will be represented by a tensor. This tensor contains one index for each of the $N/2$ even sites and the index ranges from 1 to 4, as per the even site basis. It is clear that this tensor contains the same amount of information, however in a more structured manner. The central idea of tensor networks is to represent this tensor with all of the indices as a series of contractions of smaller tensors.

\subsection{Matrix Product State}
The most common way is to construct a matrix product state (MPS) with a tensor at each lattice site with three indices, where the tensor is contracted with its two neighbors, resulting in a chain of tensors which connect all of the sites in a one-dimensional ``snake'' (see figure \ref{fig:mps}).
\begin{figure}
    \centering
    \includegraphics[width=0.6\linewidth]{figs/mps.png}
    \caption{Matrix product state with physical index $\sigma_l$ at each site\cite{PLudwig}}
    \label{fig:mps}
\end{figure}

This representation contains the same amount of physical indices $\sigma_{1,...,N}$ and can represent the same initial tensor. However, this construction is much more flexible as the dimension of the contracted indices (bond dimension) can sometimes be decreased to obtain a more compact description by eliminating redundancy. When choosing a constant bond dimension $\chi$, the amount of components now scales as $\chi^2N$ which is much better than the exponential scale seen previously.
\clearpage
\subsection{Matrix Product Operator}
Similarly to the MPS, one can construct a matrix product operator as a series of contractions of tensors for each site, whereby each site corresponds to two physical indices (see figure \ref{fig:mpo}).
\begin{figure}
    \centering
    \includegraphics[width=0.6\linewidth]{figs/mpo.png}
    \caption{Matrix product operator with physical indices $\sigma_l$ and $\sigma_l'$ at each site\cite{PLudwig}}
    \label{fig:mpo}
\end{figure}

\subsection{Example: Heisenberg Spin-1/2 System}
To see why the use of tensor networks significantly decreases computational load, consider the following example of a closed one-dimensional Heisenberg spin-$1/2$ system. The Hamiltonian is given by
\begin{equation}\label{eq:heisenberg}
    H=-\sum\limits_{i=1}^{N-1} \hat{\vec{S}}_i\cdot\hat{\vec{S}}_{i+1}=-\sum\limits_{i=1}^{N-1} \hat{S}_{x,i}\hat{S}_{x,i+1}+\hat{S}_{y,i}\hat{S}_{y,i+1}+\hat{S}_{z,i}\hat{S}_{z,i+1}
\end{equation}
using the site operators $\vec{S}_i$ which act on a two-dimensional site-local vector space.

As mentioned before, the sums in this Hamiltonian hide the fact that each summand contains many outer products with the identity. More accurately, the sum can be expanded as follows
\begin{align*}
    H=&-\mathbf{S}_{x}\otimes\mathbf{S}_{x}\otimes\mathds{1}\otimes\mathds{1}\otimes...-\mathbf{S}_{y}\otimes\mathbf{S}_{y}\otimes\mathds{1}\otimes\mathds{1}\otimes...-\mathbf{S}_{z}\otimes\mathbf{S}_{z}\otimes\mathds{1}\otimes\mathds{1}\otimes...\\
    &-\mathds{1}\otimes\mathbf{S}_{x}\otimes\mathbf{S}_{x}\otimes\mathds{1}\otimes...-\mathds{1}\otimes\mathbf{S}_{y}\otimes\mathbf{S}_{y}\otimes\mathds{1}\otimes...-\mathds{1}\otimes\mathbf{S}_{z}\otimes\mathbf{S}_{z}\otimes\mathds{1}\otimes...\\
    &-...
\end{align*}

This sum can be more efficiently represented using the following transition matrix
\begin{equation}\label{eq:transition}
    (H_\mathrm{trans})^{\alpha\beta}=\begin{pmatrix}
        1 & 0 & 0 & 0 & 0 \\
        S_x & 0 & 0 & 0 & 0 \\
        S_y & 0 & 0 & 0 & 0 \\
        S_z & 0 & 0 & 0 & 0 \\
        0 & -S_x & -S_y & -S_z & 1
    \end{pmatrix}
\end{equation}
whereby $(H_\mathrm{trans})^{\alpha\beta}_{ab}$ is actually a matrix of the site operators $(\mathbf{S}_{x,y,z})_{ab}$, so it is a rank four tensor.

Now the entire sum can be constructed as the contraction
\begin{equation}\label{eq:mpo}
    H=(H_\mathrm{trans})^{5i} (H_\mathrm{trans})^{ij}...(H_\mathrm{trans})^{kl} (H_\mathrm{trans})^{l1}
\end{equation}
where each term in the product still contains two physical indices corresponding to the site-local operator. Note that the terms at the beginning and end are row and column vectors, respectively. The equivalence of these expressions can be verified by simply multiplying out the contractions, taking into account that the multiplication of components is an outer product of site operators which does not commute.

This (exact!) representation of the Hamiltonian only requires storing $O(N)$ elements, one transition matrix for each site, which is much better than the full $O(2^N)$ components of the product basis.

\subsection{Contractions}
The previous example shows why tensor networks are apt for efficiently storing operators. However, it is not a very useful representation if it cannot be used to directly compute the necessary quantities. One important class of operations is contractions between different objects. If one naively computes the contraction while iterating over all possible sets of values of contracted indices, the exponential complexity comes back into play. Fortunately, there is a set of fairly efficient procedures which allow contracting different kinds of objects efficiently.

\subsubsection{MPS with MPS}
Firstly, the inner product of two states $|\varphi\rangle$ and $|\psi\rangle$ is a good metric for the overlap of two states. Applying the contractions in the order shown in figure \ref{fig:mps-mps}, this is done in $O(N\chi^3)$ time, which is much better than any exponential runtime in either $N$ or $\chi$.
\begin{figure}
    \centering
    \includegraphics[width=0.8\linewidth]{figs/mps-mps.png}
    \caption{Contraction procedure for the inner product of two MPSs\cite{PLudwig}}
    \label{fig:mps-mps}
\end{figure}

\subsubsection{MPO with MPS}
Next, the contraction of an operator $\hat{O}$ with a state $|\psi\rangle$ lies at the core of quantum mechanics. The procedure for this contraction is essentially the same as for the inner product, except for the second physical index of the MPO being iterated over during the contractions. This does not, however, affect the scaling of the runtime complexity. Assuming the operator has a bond dimension $W$, the runtime is now $O(NW\chi^2)$. It is important to note here that while the contraction $\hat{O}|\psi\rangle$ clearly produces another MPS efficiently, the bond dimension of the result is the product of the bond dimensions of the MPO and MPS. As such, care must be taken not to apply too many of such operations or to truncate the result, otherwise the bond dimension will grow rapidly.

\subsubsection{Expectation value of an MPO}
The expectation value of an operator $\hat{O}$ with respect to state $|\psi\rangle$ is given by the three-way contraction $\langle\psi|\hat{O}|\psi\rangle$. This operation can be viewed as a composition of the previous two and using the same procedures results in a runtime of $O(NW\chi^3)$ and $O(NW^2\chi^2)$, which is slightly worse but still not too bad.

\section{Density Matrix Renormalization Group}\label{sec:dmrg}
Density Matrix Renormalization Group (DMRG) is a procedure which is used to solve the eigenvalue problem
\begin{equation}\label{eq:dmrg}
    \hat{H}|\psi\rangle=E|\psi\rangle
\end{equation}
by solving the reduced eigenvalue problems for multiple subsystems and then combining those results. By truncating the resulting Hamiltonian based on the largest density matrix eigenvalues, the dimension of the combined system can be kept small. This procedure is then repeated until each site is included. DMRG has a very convenient implementation in the langauge of tensor networks and is implemented by the ITensor library in the Julia programming langauge\cite{ITensor}.

In order to use DMRG to solve for energy eigenstates, the MPO must be constructed for the given Hamiltonian like in the example of a Heisenberg spin system (see \ref{eq:transition}). Fortunately, an automated procedure can be applied to turn the large sum of site operators into an MPO. This is implemented by ITensor as well by means of the \verb|AutoMPO| class.

The DMRG procedure accepts many parameters which affect the precision and convergence of the solver. In addition, an initial guess is required which the solver is very sensitive to. The first parameter of importance is the number of iterations (known as `sweeps'). This parameter has a roughly linear effect on the runtime of the algorithm and, when increased, will generally improve precision. When the procedure diagonalizes the Hamiltonians of subsystems, it will do so up to a given precision. This precision parameter also has an effect on the runtime of the algorithm and the overall precision of the solution. Finally, the solver will increase or truncate the bond dimension of the solution at each sweep. Raising the maximum bond dimension will allow for more complex solutions while increasing the runtime cost. These parameters will be tuned to ensure the convergence of the algorithm.

\section{Lattice Geometry And Boundary Conditions}
Two aspects which are still ambiguous are the boundary conditions and lattice geometry. Since the MPS must cover a two-dimensional grid, the interaction distance (maximum distance between two sites which contribute to a single term in the Hamiltonian) will scale with either the lattice width or height, depending on how the MPS is constructed. The size of the bond dimension of the MPO, it turns out, roughly scales linearly with interaction distance. In order to measure the static potential as a function of distance, the height of the lattice will be varied. To obtain enough data for the fit, this dimension must span many values. As such, the lattice height will likely be quite large and should not be chosen as the layout direction for the MPS in order to reduce the bond dimension.

In addition to fixing the layout direction of the MPS, this requirement also constrains the boundary conditions along this axis to be closed. Periodic conditions would introduce the same long-range interactions again which are infeasible for simulation.

Finally, periodic boundary conditions will be chosen for the other axis in order to maximize the symmetry of the setup. Especially for small lattice widths, choosing periodic boundary conditions more closely mimics the continuum limit as translational symmetry is preserved.

A representation of the geometry and MPS layout is presented in figure \ref{fig:layout}. The main difference between the layouts used in previous work and this paper is the direct vertical paths in the lattice connecting the top and bottom. For two charges placed at equal horizontal positions at the top and bottom of the lattice, there is exactly one shortest path between the two. Another important restriction of this geometry, specifically the periodic boundary conditions, is that the width of the lattice must be even.
\begin{figure}
    \centering
    \includegraphics[width=0.4\linewidth]{figs/geom_4x6.pdf}
    \caption{Indexing scheme for MPS with axial lattice geometry. The horizontal axis wraps around to form a cylinder}
    \label{fig:layout}
\end{figure}

\section{Penalty Terms}
As mentioned in section \ref{sec:site-basis}, while the lattice sites obey Gauss' law by construction, there are still configurations where Gauss' law is violated at the odd sites in between. In order to restrict the results of DMRG to physical solutions, penalty terms will be added to the Hamiltonian for each configuration violating Gauss' law. At each odd site, the $4^3=64$ possible configurations of neighboring even lattice sites will be listed. Half of the configurations - those for which an odd number of rishons convene at the odd site - will be removed. This penalty term is constructed as follows:
\begin{align*}
    H_\mathrm{full}=H+k\sum\limits_{\text{unphysical states }|\psi\rangle} |\psi\rangle\langle\psi|
\end{align*}
where $k$ represents how strong the penalty is. This parameter must be tuned properly, as unphysical states may be allowed if it is too low and physical solutions may not be reachable if it is too high.

Referring to the pictorial representation of the lattice geometry in figure \ref{fig:layout}, it is clear that modified Gauss' law terms are needed at the closed boundaries where only one or two links meet at an odd site. These are constructed in an identical manner, except that one or two links are always set to be empty, effectively reducing the number of possible configurations to $4$ and $4^2$, respectively.

\section{Static Quark-Antiquark Potential}\label{sec:potential}
Using the methods above will yield the ground vacuum state. In order to compute the static quark-antiquark potential, the lattice must be simulated with a pair of external charges. Fortunately, these can simply be fixed at odd sites by purposely violating Gauss' law\footnote{Actually, this is not a violation of Gauss' law, since Gauss' law really states $\langle\psi|G^a_x|\psi\rangle=Q^a_x$, whereby $Q^a_x$ represents the (conserved) expected value of charge present at site $x$. This full form was omitted for clarity, but explains why this approach is valid.}. If a charge shall be placed at site $x$, all the previously unphysical configurations of neighboring states become valid and all physical configurations invalid. The penalty term will be adjusted accordingly.

In order to measure the potential energy as a function of the distance between charges, two charges will be fixed at equal horizontal positions at the top and bottom of the lattice by adjusting the relevant penalty terms. The difference between the computed energies of the empty ground state and the ground state of the system containing the two charges should provide a good estimate of the predicted potential. The vacuum state must be subtracted since the lattices have different sizes and therefore contain different amounts of vacuum energy. Note, however, that this only provides the potential as a function of lattice units, which are related to physical distance by $d=\sqrt{3}aN_y$, where $a$ represents the (unknown) physical length of a link and $N_y$ represents the height of the lattice.

As mentioned previously in section \ref{sec:dmrg}, the DMRG procedure requires a decent initial guess. For the ground state, the initial empty state provides a good estimate which respects the lattice symmetries and, more importantly, satisfies Gauss' law at all of the odd sites. This initial guess is, however, not well suited for the state containing charges, as it would violate Gauss' law near the charges. The solver will struggle to find a physical state starting from an unphysical state, as nearby states which are also unphysical may need to be traversed in order to reach a physical state. Due to the high penalty of violating Gauss' law, this is very unlikely to occur. As such, an initial guess will be chosen which connects the charges on a direct, vertical path. To relate the two simulations to one another, an initial state is chosen which represents a superposition of equal amplitude of the empty and string states.

\chapter{Results}
Results were obtained via two scripts written in Julia\footnote{Github link: \url{https://github.com/joshlengel/SU2_QLM_DMRG}. Note that this is a fork of \url{https://github.com/Liftuck/SU2_QLM_DMRG} which is the code used by P. Ludwig in his master thesis\cite{PLudwig}}. These scripts were run for different system sizes and DMRG parameters on the Marvin computing cluster maintained by the University of Bonn. Due to the shared nature of cluster, computational resources are limited. Therefore, individual jobs are limited to run for seven days.

\section{Lattice Sizing and Coupling Strength}
Nevertheless, a lattice width of $N_x=4$ was chosen and simulations were run for $3\leq N_y\leq 12$. Simulations with $N_x=6$ were attempted but only $3\leq N_y\leq 6$ actually completed in a reasonable amount of time. This is not enough data to analyze so only the $N_x=4$ cases will be studied. Each lattice size was simulated for 31 different values of $g^2$, ranging linearly between 0.5 and 8, inclusive.

\section{DMRG Parameters}
Each simulation runs for 10000 sweeps (iterations) but terminates early if the energy of the state changes by less than $5\cdot 10^{-5}$ from the previous sweep and a minimum number of sweeps have been performed. For each lattice size, the simulation is run for the highest $g^2$ with the initial state specified in section \ref{sec:potential} and for a minimum of 51 sweeps. Additionally, random noise is applied to the result of the first 15 sweeps to avoid the solver getting stuck in local minima. Then, the lattice is simulated for each remaining $g^2$ in descending order using the previous solution as the initial state. The minimum number of sweeps is reduced to 7 and only 4 noise sweeps are performed.

$\left\lceil \frac{5}{4}N_xN_y\right\rceil$ was found previously to be a good value for the initial maximum link dimension. Furthermore, the maximum link dimension is rapidly increased from sweep to sweep. This scheme is taken in its entirety from P. Ludwig's master thesis\cite{PLudwig}.

The simulations were carried out with a cutoff (precision) of $10^{-9}$ at first, however the lattices with $3\leq N_y\leq 8$ were simulated again with a lower cutoff of $10^{-12}$.

\section{Exact Diagonalization}
In order to verify the correctness of the procedure, the vacuum Hamiltonian was diagonalized exactly using \verb|scipy.sparse.linalg.eigsh| from the Python library SciPy\cite{SciPy}. Firstly, all $4^{16}=2^{32}$ basis states are listed. Next, all basis states violating Gauss' law are removed from the list, resulting in $2^{13}$ states remaining. Then, since each basis state is already an eigenstate of the electric Hamiltonian, the number of active links is counted and multiplied by the flux of a single link, $2\cdot 6/4$, and finally entered into the corresponding entry on the diagonal.

The magnetic Hamiltonian is generated by considering the 64x64 matrix constructed previously for a single plaquette. Then, each of the 12 plaquettes in the 4x4 system is iterated over and an entry added for each non-zero transition in the plaquette Hamiltonian. Finally, the two matrices are weighted by $g^2$ and $1/g^2$ and summed.

The ground state and mass gap are presented in figure \ref{fig:massgap} as a function of $g^2$.
\begin{figure}
    \centering
    \includegraphics[width=0.95\linewidth]{figs/massgaps.png}
    \caption{$g^2$-dependent ground state energy (left) and mass gap (right) of 4x4 lattice using exact diagonalization}
    \label{fig:massgap}
\end{figure}

\section{Verification}
Before the results can be analyzed, the solutions found by DMRG must be checked for validity. Since DMRG should return the ground state, consider the solution
$$|\psi\rangle=\sqrt{f}|\psi_0\rangle+\sqrt{1-f}|\psi_1\rangle$$
which represents the superposition of the actual ground state $|\psi_0\rangle$ and the first excited state $|\psi_1\rangle$ with the weight factor $f$, known as the fidelity. It follows that
\begin{align*}
    \mathrm{Var}\left[\hat{H}\right]&=\langle\psi|\hat{H}^2|\psi\rangle-\left(\langle\psi|\hat{H}|\psi\rangle\right)^2=f{E_0}^2+(1-f){E_1}^2-\left(fE_0+(1-f)E_1\right)^2\\
    &=f(1-f){E_0}^2+f(1-f){E_1}^2-2f(1-f)E_0E_1=f(1-f)\left(E_1-E_0\right)^2\\
    &=(1-f)\left(E_1-E_0\right)^2-(1-f)^2\left(E_1-E_0\right)^2\\
    &\approx (1-f)\left(E_1-E_0\right)^2\\
    \implies 1-f&\approx \frac{\mathrm{Var}\left[\hat{H}\right]}{\left(E_1-E_0\right)^2}\equiv \frac{\mathrm{Var}\left[\hat{H}\right]}{\Delta m^2}
\end{align*}
assuming the state is reasonably well converged ($1-f\ll 1$). The true mass gap $E_1-E_0$ is not known, but the values obtained by exact diagonalization of the 4x4 lattice can be used as an approximation. A well converged state must satisfy $1-f\ll 1$ in order to be considered as a solution. As figure \ref{fig:fidelity} clearly shows, the fidelity is reasonably high for each solution. There is little value in estimating the statistical error of DMRG in this way since the mass gaps used are approximate for $N_y\leq 4$ and the fidelity does not account for finite-size effects.
\begin{figure}[htbp]
    \centering
    \includegraphics[width=0.95\linewidth]{figs/fidelity.png}
    \caption{$1-f$ of each vacuum (left) and string (right) solution}
    \label{fig:fidelity}
\end{figure}

Additionally, a state must satisfy Gauss' law at the odd sites to be considered physical. This is enforced via the penalty terms, however there is no guarantee that any solution of the DMRG actually obeys Gauss' law. As such, the expectation value of the penalty terms will be checked to ensure the validity of the solution. Once again, figure \ref{fig:gauss} clearly shows that Gauss' law holds for all solutions. Specifically, the ground state systems are subject to an error of below $10^{-12}$ for $3<N_y\leq 6$ and $10^{-9}$ for the remaining lattices, in correspondence with the precision of the DMRG. For the external charges, the penalty factor of Gauss' law is 10000 times higher than the others in order to ensure the string is not broken. Due to the slower convergence of the string states, the error is much higher around $10^{-8}$. This is still minuscule in comparison to the electric fluxes being measured, so all solutions are valid.
\begin{figure}
    \centering
    \includegraphics[width=0.95\linewidth]{figs/gauss.png}
    \caption{Expectation value of penalty operator of each vacuum (left) and string (right) solution}
    \label{fig:gauss}
\end{figure}

\section{Potentials}
Figure \ref{fig:potential} displays all of the computed potentials.
\begin{figure}
    \centering
    \includegraphics[width=0.6\linewidth]{figs/potential.png}
    \caption{Potential $\Delta E$ for different values of $g^2$ and $N_y$}
    \label{fig:potential}
\end{figure}

The potential is minimal at $4.75\leq g^2\leq 5.0$ for a fixed lattice height. Near $g^2=0$ the magnetic term dominates and a $1/g^2$ dependence is observed. On the contrary, for large $g^2$, the electric term dominates and since it is already diagonal, the energy is simply $2\cdot 6/4\cdot N_y g^2$. These eigenstates are not degenerate due to the fact that there is only a single shortest path between the two charges.

Furthermore, the potential increases with increasing $N_y$ for all values of $g^2$. This is to be expected with the Cornell potential, as long as the string tension $\sigma$ is positive and the L\"uscher term $\gamma$ is either negative or too small to be observed.

Figure \ref{fig:error} shows that the difference between the ground state energy of the vacuum computed by DMRG and by exact diagonalization is on the order of $10^{-9}$. The positive linear correlation is unexpected and opposes intuition, since the energies are higher at lower $g^2$ values and DMRG converges much more quickly for high $g^2$ values. This is explained by observing that the relative error is on the order of $10^{-11}$ or $10^{-12}$, which is the same order as the precision of the DMRG. In conjunction with the precision of the eigensolver from SciPy, these errors lie below the precision expected of either procedure, thus affirming the validity of the results.
\begin{figure}
    \centering
    \includegraphics[width=0.6\linewidth]{figs/error.png}
    \caption{Difference of $g^2$-dependent vacuum ground state energies of 4x4 lattice using DMRG and exact diagonalization.}
    \label{fig:error}
\end{figure}

\section{L{\"u}scher Term}
Next, for each $g^2$, the potential is written as a function of $N_y$. Afterwards, the Cornell potential is fitted to the data using \verb|scipy.optimize.curve_fit| from the Python library SciPy\cite{SciPy} to determine the values for $\gamma$, $\sigma$ and $c$. Since DMRG does not provide an estimate for the statistical uncertainty of the eigenvalues, a relative error $\Delta E/E$ is chosen which results in a reduced $\chi^2$ of one. One such potential with $g^2=5.0$ is presented in Figure \ref{fig:cornell-fit}, whereby the bottom-most graph subtracts the linear component, thereby isolating the L\"uscher term.
\begin{figure}[htbp]
    \centering
    \includegraphics[width=0.95\linewidth]{figs/cornell-fit.png}
    \includegraphics[width=0.475\linewidth]{figs/cornell-luescher.png}
    \caption{Potential as a function of $N_y$ for $g^2=5.0$ (top left), residues of data w.r.t. the fitting function (top right) and potential with subtracted linear component $\sigma N_y+c$ (bottom center)}
    \label{fig:cornell-fit}
\end{figure}

This fit is repeated for each value of $g^2$ and the resulting parameters plotted as functions of $g^2$ (see figure \ref{fig:cornell-params}).
\begin{figure}[htbp]
    \centering
    \includegraphics[width=0.9\linewidth]{figs/gamma.png}
    \includegraphics[width=0.9\linewidth]{figs/sigma-c.png}
    \caption{Parameters $\gamma$, $\sigma$ and $c$ as functions of $g^2$}
    \label{fig:cornell-params}
\end{figure}

The $N_y=3$ value was excluded in each of the fits on the grounds of finite size effects. As expected, the L\"uscher term is negative and the string tension is positive. The high errors of for $\gamma$ (estimated by the fitting routine) indicate a low confidence in this result. In fact, if one excludes the $N_y=4$ data point from the fits, the behaviour of $\gamma$ is flipped (see figures \ref{fig:cornell-fit-excluded} and \ref{fig:cornell-params-excluded}).

Since the L\"uscher term is only relevant for small lattice heights, one expects $\gamma$ to go to zero as more small lattices are excluded. The fact that $\gamma$ has become positive, however with only approximately half of the amplitude, implies that the form of the graph observed is merely a consequence of statistical error and not of any significance.
\begin{figure}[htbp]
    \centering
    \includegraphics[width=0.95\linewidth]{figs/cornell-fit-excluded.png}
    \includegraphics[width=0.475\linewidth]{figs/cornell-luescher-excluded.png}
    \caption{Potential as a function of $N_y$ for $g^2=5.0$ (top left), residues of data w.r.t. the fitting function (top right) and potential with subtracted linear component $\sigma N_y+c$ (bottom center) excluding $N_y=4$ data point}
    \label{fig:cornell-fit-excluded}
\end{figure}
\clearpage

Another issue with the observed values of $\gamma$ is the fact that they do not appear to match the prediction of $\gamma=-\frac{\pi}{24}$ (or $-\frac{\pi}{12}$, as argued by Di Marcantonio\cite{DiMarcantonio}) for small coupling strengths.
\begin{figure}[htbp]
    \centering
    \includegraphics[width=0.9\linewidth]{figs/gamma-excluded.png}
    \includegraphics[width=0.9\linewidth]{figs/sigma-c-excluded.png}
    \caption{Parameters $\gamma$, $\sigma$ and $c$ as functions of $g^2$ excluding $N_y=4$ data point}
    \label{fig:cornell-params-excluded}
\end{figure}

Overall, it seems that a L\"uscher term is not present in the data. The data does, however, not exclude the possibility of a L\"uscher term. It appears that larger lattice widths may be needed to reduce finite size effects. Furthermore, the lattice height could be kept constant and the charges placed within the lattice. This approach may mitigate the effects which stem from the different lattice sizes.

\section{String Width}
Having excluded the existence of a L\"uscher term in the data, the shape of the strings will be analyzed to determine whether the system possesses any phase structure which may explain the absence of the L\"uscher term.

\subsection{Cross-section}
Consider a horizontal slice centered vertically as in figure \ref{fig:cross-section}. The expectation value of the operator $|i\rangle\langle i|+|j\rangle\langle j|$ is computed at each site near the slice, such that $i$ and $j$ represent basis states where the link intersecting the cross section is active. This effectively computes the expected number of active links in the cross section at each link.
\begin{figure}
    \centering
    \includegraphics[width=0.5\linewidth]{figs/cross-section.pdf}
    \caption{Cross section (gray) for measuring string width}
    \label{fig:cross-section}
\end{figure}

Due to the geometry of the cross section, only four unique data points can be obtained per slice. To improve the results, the link opposite the string is duplicated in order to symmetrize the data (this is valid due to the periodic boundary conditions). The Gaussian function
$$f(x)=Ae^{-\frac{x^2}{2\omega^2}}+c$$
is fitted to the cross section for each lattice size and coupling strength. $x$ represents the horizontal distance from the string, whereby the distance between two neighboring horizontal links is $1.5$ lattice units. $A$ represents the maximum amplitude of the string which is limited to the range $[0,1]$, $\omega^2$ is the string width and $c$ is a background which arises from the small finite extent of the transverse dimension. This function is generally a bad model however, since it is not smooth across the periodic boundary (the derivative is discontinuous on the opposite side of the string). As such, the following ``compact'' Gaussian will be fitted as well:
$$f(x)=\left\{\begin{array}{cc}Ae^{-\frac{2x^2}{4\omega^2-x^2}}+c & |x|<2\omega \\ 0 & |x|\geq 2\omega \end{array}\right.$$
This function reduces to the standard Gaussian to second order in $x/\omega$, but compresses the infinite extent of the Gaussian into the range $[-2\omega,2\omega]$, outside of which it is exactly zero. As such, assuming the extent of $4\omega$ is smaller than the lattice width, this function is smooth and periodic.

It should be noted that the expectation values being measured along the cross section do not have an intrinsic statistical error. As such, no $\chi^2$ can be computed for any fits done. While fixing $\chi^2=1$ and using this condition to set a relative error would mitigate this issue, this procedure does not provide much value here since the low number of data points restricts this analysis to the quantitative realm anyway.

\subsection{Rough Strings}
The strings can either be rigid or rough, whereby the width of rough strings scale as $\omega^2=A\ln(N_y/b)$, whereby $A$ and $b$ are proportionality parameters. Importantly, a non-zero L\"uscher term generally implies the existence of rough strings in the range $g_\mathrm{c}\leq g\leq g_\mathrm{r}$, with $g_\mathrm{r}$ denoting the roughening transition and $g_\mathrm{c}$ representing the phase transition to a deconfined model\cite{DiMarcantonio}. As such, the dependence of the string width on the lattice height could give valuable insight on whether a L\"uscher term could be found using this setup with improved precision.

For brevity, graphical depictions of the cross section for each $g^2$ and $N_y$ are relegated to the Appendix. Figure \ref{fig:string-width} shows the string width as a function of lattice height for $g^2=2$. While the data does not have any obvious outliers, the data points oscillate between even and odd lattice heights. This is due to the fact that the vertical center of the string only moves by half of a lattice unit for every unit the lattice grows. As such, the cross sections pass through the string at different angles for even and odd heights. Considering also the poor fitting results for small lattice heights, only odd lattice heights larger than four will be analysed in the following.
\begin{figure}[htbp]
    \centering
    \includegraphics[width=0.95\linewidth]{figs/string-width.png}
    \caption{String width computed using Gaussian (left) and compact Gaussian (right) as a function of string length (lattice height) for $g^2=2$}
    \label{fig:string-width}
\end{figure}

Figure \ref{fig:omega} shows the computed string widths using the Gaussian and compact Gaussian for multiple values of $g^2$ as a function of the string length (lattice height).
\begin{figure}[htbp]
    \centering
    \includegraphics[width=0.95\linewidth]{figs/omega.png}
    \caption{String width computed using Gaussian (left) and compact Gaussian (right) as a function of string length (lattice height)}
    \label{fig:omega}
\end{figure}
\clearpage

For low coupling strengths, the data appears to fit the logarithmic model. Noticeably, the compact Gaussian consistently estimates larger string widths than the standard Gaussian. This makes sense, as the compact Gaussian falls to zero much quicker thereby allowing for a wider string width near the center. Moreover, this shows the shortcomings of the lack of data, as many different models could conceivably fit this data with very different results. Figure \ref{fig:omega-diff} displays the (signed) difference in string width $\omega^2$ as determined by the compact and standard Gaussian functions. As expected, this difference is positive for all values, however there appears to be a minimum around $g^2=6.25$ which coincides with a poor fit of the rough string model for $5.5\leq g^2\leq 6.5$.
\begin{figure}[htbp]
    \centering
    \includegraphics[width=0.5\linewidth]{figs/omega-diff.png}
    \caption{Difference of $\omega^2$ as determined by compact and standard Gaussian functions}
    \label{fig:omega-diff}
\end{figure}

\subsection{Phase Structure}
As shown by figure \ref{fig:string-params}, there appears to be lack of convergence of the rough string model for $5.0\leq g^2\leq 6.5$. Interestingly, the compact Gaussian model is able to converge for larger $g^2$ values while the standard Gaussian is not. This is due to the fact that only the string center is measurably active while all other data points are, for all intents and purposes, zero. The string width of the standard Gaussian is therefore very sensitive to small deviations from zero, while the compact Gaussian measures a mostly consistent string width of $1.5/2$ lattice units. In both cases, the values should not be trusted, however, as the string is too thin to reliably determine the string width. While this could be interpreted as the roughening transition (oscillation width dropping below a lattice unit), it is likely this behaviour would change somewhat with increased lattice width. This does, however, provide a probable upper bound on the roughening transition $g_\mathrm{r}^2<6.5\implies g_\mathrm{r}<2.55$.
\begin{figure}[htbp]
    \centering
    \includegraphics[width=0.95\linewidth]{figs/string-params.png}
    \caption{Rough string parameters $A$ and $b$ as functions of $g^2$}
    \label{fig:string-params}
\end{figure}

Below the roughening transition, both models agree on the proportionality constant $A$. Taking the difference of the squared string widths yields
\begin{align*}
    \omega_\mathrm{compact}^2-\omega_\mathrm{gauss}^2&=(A_\mathrm{compact}-A_\mathrm{gauss})\ln N_y - A_\mathrm{compact}\ln b_\mathrm{compact} + A_\mathrm{gauss}\ln b_\mathrm{gauss}\\
    &=\Delta A\ln N_y + \bar{A}\ln\left(b_\mathrm{gauss}/b_\mathrm{compact}\right)-\frac{1}{2}\Delta A\ln\left(b_\mathrm{gauss}b_\mathrm{compact}\right)
\end{align*}
whereby $\Delta A\equiv A_\mathrm{compact}-A_\mathrm{gauss}$ and $\bar{A}\equiv \frac{1}{2}\left(A_\mathrm{compact}+A_\mathrm{gauss}\right)$. If the strings are rough, both Gaussians should generally produce the same functional dependence between $N_y$ and $\omega^2$. As such, one expects the difference $\omega_\mathrm{compact}^2-\omega_\mathrm{gauss}^2$ not to depend on $N_y$, which implies $\Delta A\approx 0$ as a necessary condition. Using this fact, only the second term remains resulting in $\Delta(\omega^2)\approx \bar{A}\ln\left(b_\mathrm{gauss}/b_\mathrm{compact}\right)$, which decidedly depends neither on $\omega$, $N_y$ or $g^2$.

Since only five data points are available for each fit, one expects the outermost points to lie at roughly zero, the middle point to be at the maximum and for the two remaining to determine the value of $\omega$. Comparing the two Gaussians in this manner reveals
$$e^{-\frac{x^2}{2\omega_\mathrm{gauss}^2}}=e^{-\frac{2x^2}{4\omega_\mathrm{compact}^2-x^2}}\implies \omega_\mathrm{compact}^2-\omega_\mathrm{gauss}^2=\frac{x^2}{4}=\frac{1.5^2}{4}=0.5625$$

Figure \ref{fig:string-params-diff} shows that in the roughening region, the first condition holds within a 10\% margin, with $\Delta A/A<1\%$ for $g^2>3.5$. Furthermore, $\Delta(\omega^2)$ stays roughy constant in the roughening region, only varying by approximately $22\%$ from the predicted value.
\begin{figure}
    \centering
    \includegraphics[width=0.95\linewidth]{figs/string-params-diff.png}
    \caption{Comparison of rough string parameters for standard and compact Gaussian}
    \label{fig:string-params-diff}
\end{figure}

This suggests that the strings do follow the logarithmic scaling independently of how the string width is specifically defined and that the strings are, in fact, rough strings. It follows that there should be a L\"uscher term, which may however be overshadowed by finite size effects.

\chapter{Conclusion}
Firstly, the lattice sizing and DMRG parameters were chosen to tune the balance between lattice size and computational resources required. A lattice width of $N_x=4$ was found to be the largest width which could feasibly be simulated with the available resources. After multiple weeks of simulations with a DMRG cutoff of $10^{-9}$ the vacuum and string states of lattices with heights $3\leq N_y\leq 12$ were found for 31 linearly spaced coupling strengths $0.5\leq g^2\leq 8.0$. The lattices with $4\leq N_y\leq 8$ were simulated again with a lower cutoff of $10^{-12}$ for increased precision. All simulations converged nicely with vanishing penalty terms and reasonably high fidelity.

Plotting the computed Cornell potential - the difference between the energies of the string state and vacuum state as a function of the charge separation $r$ - it was evident that no L\"uscher term ($\propto \frac{1}{r}$) could be extracted with confidence, as excluding one data point resulted in a sign flip and the statistical errors were substantial. Finite size effects appeared to be a relevant factor in these results which presents an opportunity for further research into this lattice geometry. Furthermore, fixing the lattice height and placing the charges at varying distances within the lattice may prove to deliver more definitive results. Another aspect to consider is the symmetry of the lattice. Strings are no longer symmetric about the vertical line connecting the two charges, which suggests that other string placements could provide more data.

Finally, the string width was determined using a horizontal cross section of the string by determining the expectation value of the number of active links passing through the cross section at each link. A standard Gaussian resulted in fits with very volatile parameters, so a ``compact'' Gaussian was used instead with a finite non-zero region which has the added benefit of being periodic. The resulting string widths differed by a substantial amount, but seemed to coincide in the region $6.0\leq g^2\leq 6.5$, after which the data became too narrow for the models to be deemed accurate.

For $g^2\leq 6.5$, the string width qualitatively followed a logarithmic model with respect to the string length, which suggests that a L\"uscher term should exist here. after this point, the models broke down, suggesting that the strings became rigid. This coincides with an upper bound of $g_\mathrm{r}\leq 2.55$ for the roughening transition, which is loosely in agreement with $g_\mathrm{r}\approx 1$ found by previous studies\cite{DiMarcantonio}.

In summary, the existence of an axial L\"uscher term is inconclusive for a hexagonal lattice of width four. Nevertheless, the analysis of the string width suggests that a L\"uscher term should exist and that it may be found using other methods, which opens the door for further research into SU(2) quantum link models on a hexagonal lattice.

\printbibliography

\chapter*{Appendix}
\begin{figure}[!hb]
    \centering
    \includegraphics[width=\linewidth]{figs/V_g2_0.5.png}
\end{figure}
\begin{figure}[!hb]\ContinuedFloat
    \centering
    \includegraphics[width=\linewidth]{figs/V_g2_0.75.png}
\end{figure}
\begin{figure}[!hb]\ContinuedFloat
    \centering
    \includegraphics[width=\linewidth]{figs/V_g2_1.0.png}
\end{figure}
\begin{figure}[!hb]\ContinuedFloat
    \centering
    \includegraphics[width=\linewidth]{figs/V_g2_1.25.png}
\end{figure}
\begin{figure}[!hb]\ContinuedFloat
    \centering
    \includegraphics[width=\linewidth]{figs/V_g2_1.5.png}
\end{figure}
\begin{figure}[!hb]\ContinuedFloat
    \centering
    \includegraphics[width=\linewidth]{figs/V_g2_1.75.png}
\end{figure}
\begin{figure}[!hb]\ContinuedFloat
    \centering
    \includegraphics[width=\linewidth]{figs/V_g2_2.0.png}
\end{figure}
\begin{figure}[!hb]\ContinuedFloat
    \centering
    \includegraphics[width=\linewidth]{figs/V_g2_2.25.png}
\end{figure}
\begin{figure}[!hb]\ContinuedFloat
    \centering
    \includegraphics[width=\linewidth]{figs/V_g2_2.5.png}
\end{figure}
\begin{figure}[!hb]\ContinuedFloat
    \centering
    \includegraphics[width=\linewidth]{figs/V_g2_2.75.png}
\end{figure}
\begin{figure}[!hb]\ContinuedFloat
    \centering
    \includegraphics[width=\linewidth]{figs/V_g2_3.0.png}
\end{figure}
\begin{figure}[!hb]\ContinuedFloat
    \centering
    \includegraphics[width=\linewidth]{figs/V_g2_3.25.png}
\end{figure}
\begin{figure}[!hb]\ContinuedFloat
    \centering
    \includegraphics[width=\linewidth]{figs/V_g2_3.5.png}
\end{figure}
\begin{figure}[!hb]\ContinuedFloat
    \centering
    \includegraphics[width=\linewidth]{figs/V_g2_3.75.png}
\end{figure}
\begin{figure}[!hb]\ContinuedFloat
    \centering
    \includegraphics[width=\linewidth]{figs/V_g2_4.0.png}
\end{figure}
\begin{figure}[!hb]\ContinuedFloat
    \centering
    \includegraphics[width=\linewidth]{figs/V_g2_4.25.png}
\end{figure}
\begin{figure}[!hb]\ContinuedFloat
    \centering
    \includegraphics[width=\linewidth]{figs/V_g2_4.5.png}
\end{figure}
\begin{figure}[!hb]\ContinuedFloat
    \centering
    \includegraphics[width=\linewidth]{figs/V_g2_4.75.png}
\end{figure}
\begin{figure}[!hb]\ContinuedFloat
    \centering
    \includegraphics[width=\linewidth]{figs/V_g2_5.0.png}
\end{figure}
\begin{figure}[!hb]\ContinuedFloat
    \centering
    \includegraphics[width=\linewidth]{figs/V_g2_5.25.png}
\end{figure}
\begin{figure}[!hb]\ContinuedFloat
    \centering
    \includegraphics[width=\linewidth]{figs/V_g2_5.5.png}
\end{figure}
\begin{figure}[!hb]\ContinuedFloat
    \centering
    \includegraphics[width=\linewidth]{figs/V_g2_5.75.png}
\end{figure}
\begin{figure}[!hb]\ContinuedFloat
    \centering
    \includegraphics[width=\linewidth]{figs/V_g2_6.0.png}
\end{figure}
\begin{figure}[!hb]\ContinuedFloat
    \centering
    \includegraphics[width=\linewidth]{figs/V_g2_6.25.png}
\end{figure}
\begin{figure}[!hb]\ContinuedFloat
    \centering
    \includegraphics[width=\linewidth]{figs/V_g2_6.5.png}
\end{figure}
\begin{figure}[!hb]\ContinuedFloat
    \centering
    \includegraphics[width=\linewidth]{figs/V_g2_6.75.png}
\end{figure}
\begin{figure}[!hb]\ContinuedFloat
    \centering
    \includegraphics[width=\linewidth]{figs/V_g2_7.0.png}
\end{figure}
\begin{figure}[!hb]\ContinuedFloat
    \centering
    \includegraphics[width=\linewidth]{figs/V_g2_7.25.png}
\end{figure}
\begin{figure}[!hb]\ContinuedFloat
    \centering
    \includegraphics[width=\linewidth]{figs/V_g2_7.5.png}
\end{figure}
\begin{figure}[!hb]\ContinuedFloat
    \centering
    \includegraphics[width=\linewidth]{figs/V_g2_7.75.png}
\end{figure}
\begin{figure}[!hb]\ContinuedFloat
    \centering
    \includegraphics[width=\linewidth]{figs/V_g2_8.0.png}
    \caption{Potential energy as a function of $N_y$ for different values of $g^2$}
\end{figure}

\begin{figure}[!hb]
    \centering
    \begin{minipage}{\linewidth}
        \includegraphics[width=0.32\linewidth]{figs/s_g2_0.5_Ny_3.png}
        \includegraphics[width=0.32\linewidth]{figs/s_g2_0.75_Ny_3.png}
        \includegraphics[width=0.32\linewidth]{figs/s_g2_1.0_Ny_3.png}
    \end{minipage}
\end{figure}
\begin{figure}[!hb]\ContinuedFloat
    \centering
    \begin{minipage}{\linewidth}
        \includegraphics[width=0.32\linewidth]{figs/s_g2_0.5_Ny_4.png}
        \includegraphics[width=0.32\linewidth]{figs/s_g2_0.75_Ny_4.png}
        \includegraphics[width=0.32\linewidth]{figs/s_g2_1.0_Ny_4.png}
    \end{minipage}
\end{figure}
\begin{figure}[!hb]\ContinuedFloat
    \centering
    \begin{minipage}{\linewidth}
        \includegraphics[width=0.32\linewidth]{figs/s_g2_0.5_Ny_5.png}
        \includegraphics[width=0.32\linewidth]{figs/s_g2_0.75_Ny_5.png}
        \includegraphics[width=0.32\linewidth]{figs/s_g2_1.0_Ny_5.png}
    \end{minipage}
\end{figure}
\begin{figure}[!hb]\ContinuedFloat
    \centering
    \begin{minipage}{\linewidth}
        \includegraphics[width=0.32\linewidth]{figs/s_g2_0.5_Ny_6.png}
        \includegraphics[width=0.32\linewidth]{figs/s_g2_0.75_Ny_6.png}
        \includegraphics[width=0.32\linewidth]{figs/s_g2_1.0_Ny_6.png}
    \end{minipage}
\end{figure}
\begin{figure}[!hb]\ContinuedFloat
    \centering
    \begin{minipage}{\linewidth}
        \includegraphics[width=0.32\linewidth]{figs/s_g2_0.5_Ny_7.png}
        \includegraphics[width=0.32\linewidth]{figs/s_g2_0.75_Ny_7.png}
        \includegraphics[width=0.32\linewidth]{figs/s_g2_1.0_Ny_7.png}
    \end{minipage}
\end{figure}
\begin{figure}[!hb]\ContinuedFloat
    \centering
    \begin{minipage}{\linewidth}
        \includegraphics[width=0.32\linewidth]{figs/s_g2_0.5_Ny_8.png}
        \includegraphics[width=0.32\linewidth]{figs/s_g2_0.75_Ny_8.png}
        \includegraphics[width=0.32\linewidth]{figs/s_g2_1.0_Ny_8.png}
    \end{minipage}
\end{figure}
\begin{figure}[!hb]\ContinuedFloat
    \centering
    \begin{minipage}{\linewidth}
        \includegraphics[width=0.32\linewidth]{figs/s_g2_0.5_Ny_9.png}
        \includegraphics[width=0.32\linewidth]{figs/s_g2_0.75_Ny_9.png}
        \includegraphics[width=0.32\linewidth]{figs/s_g2_1.0_Ny_9.png}
    \end{minipage}
\end{figure}
\begin{figure}[!hb]\ContinuedFloat
    \centering
    \begin{minipage}{\linewidth}
        \includegraphics[width=0.32\linewidth]{figs/s_g2_0.5_Ny_10.png}
        \includegraphics[width=0.32\linewidth]{figs/s_g2_0.75_Ny_10.png}
        \includegraphics[width=0.32\linewidth]{figs/s_g2_1.0_Ny_10.png}
    \end{minipage}
\end{figure}
\begin{figure}[!hb]\ContinuedFloat
    \centering
    \begin{minipage}{\linewidth}
        \includegraphics[width=0.32\linewidth]{figs/s_g2_0.5_Ny_11.png}
        \includegraphics[width=0.32\linewidth]{figs/s_g2_0.75_Ny_11.png}
        \includegraphics[width=0.32\linewidth]{figs/s_g2_1.0_Ny_11.png}
    \end{minipage}
\end{figure}
\begin{figure}[!hb]\ContinuedFloat
    \centering
    \begin{minipage}{\linewidth}
        \includegraphics[width=0.32\linewidth]{figs/s_g2_0.5_Ny_12.png}
        \includegraphics[width=0.32\linewidth]{figs/s_g2_0.75_Ny_12.png}
        \includegraphics[width=0.32\linewidth]{figs/s_g2_1.0_Ny_12.png}
    \end{minipage}
    \caption{String cross section for $g^2=0.5,\,0.75,\,1.0$ (standard Gaussian in blue, compact Gaussian in orange)}
\end{figure}

\begin{figure}[!hb]
    \centering
    \begin{minipage}{\linewidth}
        \includegraphics[width=0.32\linewidth]{figs/s_g2_1.25_Ny_3.png}
        \includegraphics[width=0.32\linewidth]{figs/s_g2_1.5_Ny_3.png}
        \includegraphics[width=0.32\linewidth]{figs/s_g2_1.75_Ny_3.png}
    \end{minipage}
\end{figure}
\begin{figure}[!hb]\ContinuedFloat
    \centering
    \begin{minipage}{\linewidth}
        \includegraphics[width=0.32\linewidth]{figs/s_g2_1.25_Ny_4.png}
        \includegraphics[width=0.32\linewidth]{figs/s_g2_1.5_Ny_4.png}
        \includegraphics[width=0.32\linewidth]{figs/s_g2_1.75_Ny_4.png}
    \end{minipage}
\end{figure}
\begin{figure}[!hb]\ContinuedFloat
    \centering
    \begin{minipage}{\linewidth}
        \includegraphics[width=0.32\linewidth]{figs/s_g2_1.25_Ny_5.png}
        \includegraphics[width=0.32\linewidth]{figs/s_g2_1.5_Ny_5.png}
        \includegraphics[width=0.32\linewidth]{figs/s_g2_1.75_Ny_5.png}
    \end{minipage}
\end{figure}
\begin{figure}[!hb]\ContinuedFloat
    \centering
    \begin{minipage}{\linewidth}
        \includegraphics[width=0.32\linewidth]{figs/s_g2_1.25_Ny_6.png}
        \includegraphics[width=0.32\linewidth]{figs/s_g2_1.5_Ny_6.png}
        \includegraphics[width=0.32\linewidth]{figs/s_g2_1.75_Ny_6.png}
    \end{minipage}
\end{figure}
\begin{figure}[!hb]\ContinuedFloat
    \centering
    \begin{minipage}{\linewidth}
        \includegraphics[width=0.32\linewidth]{figs/s_g2_1.25_Ny_7.png}
        \includegraphics[width=0.32\linewidth]{figs/s_g2_1.5_Ny_7.png}
        \includegraphics[width=0.32\linewidth]{figs/s_g2_1.75_Ny_7.png}
    \end{minipage}
\end{figure}
\begin{figure}[!hb]\ContinuedFloat
    \centering
    \begin{minipage}{\linewidth}
        \includegraphics[width=0.32\linewidth]{figs/s_g2_1.25_Ny_8.png}
        \includegraphics[width=0.32\linewidth]{figs/s_g2_1.5_Ny_8.png}
        \includegraphics[width=0.32\linewidth]{figs/s_g2_1.75_Ny_8.png}
    \end{minipage}
\end{figure}
\begin{figure}[!hb]\ContinuedFloat
    \centering
    \begin{minipage}{\linewidth}
        \includegraphics[width=0.32\linewidth]{figs/s_g2_1.25_Ny_9.png}
        \includegraphics[width=0.32\linewidth]{figs/s_g2_1.5_Ny_9.png}
        \includegraphics[width=0.32\linewidth]{figs/s_g2_1.75_Ny_9.png}
    \end{minipage}
\end{figure}
\begin{figure}[!hb]\ContinuedFloat
    \centering
    \begin{minipage}{\linewidth}
        \includegraphics[width=0.32\linewidth]{figs/s_g2_1.25_Ny_10.png}
        \includegraphics[width=0.32\linewidth]{figs/s_g2_1.5_Ny_10.png}
        \includegraphics[width=0.32\linewidth]{figs/s_g2_1.75_Ny_10.png}
    \end{minipage}
\end{figure}
\begin{figure}[!hb]\ContinuedFloat
    \centering
    \begin{minipage}{\linewidth}
        \includegraphics[width=0.32\linewidth]{figs/s_g2_1.25_Ny_11.png}
        \includegraphics[width=0.32\linewidth]{figs/s_g2_1.5_Ny_11.png}
        \includegraphics[width=0.32\linewidth]{figs/s_g2_1.75_Ny_11.png}
    \end{minipage}
\end{figure}
\begin{figure}[!hb]\ContinuedFloat
    \centering
    \begin{minipage}{\linewidth}
        \includegraphics[width=0.32\linewidth]{figs/s_g2_1.25_Ny_12.png}
        \includegraphics[width=0.32\linewidth]{figs/s_g2_1.5_Ny_12.png}
        \includegraphics[width=0.32\linewidth]{figs/s_g2_1.75_Ny_12.png}
    \end{minipage}
    \caption{String cross section for $g^2=1.25,\,1.5,\,1.75$ (standard Gaussian in blue, compact Gaussian in orange)}
\end{figure}

\begin{figure}[!hb]
    \centering
    \begin{minipage}{\linewidth}
        \includegraphics[width=0.32\linewidth]{figs/s_g2_2.0_Ny_3.png}
        \includegraphics[width=0.32\linewidth]{figs/s_g2_2.25_Ny_3.png}
        \includegraphics[width=0.32\linewidth]{figs/s_g2_2.5_Ny_3.png}
    \end{minipage}
\end{figure}
\begin{figure}[!hb]\ContinuedFloat
    \centering
    \begin{minipage}{\linewidth}
        \includegraphics[width=0.32\linewidth]{figs/s_g2_2.0_Ny_4.png}
        \includegraphics[width=0.32\linewidth]{figs/s_g2_2.25_Ny_4.png}
        \includegraphics[width=0.32\linewidth]{figs/s_g2_2.5_Ny_4.png}
    \end{minipage}
\end{figure}
\begin{figure}[!hb]\ContinuedFloat
    \centering
    \begin{minipage}{\linewidth}
        \includegraphics[width=0.32\linewidth]{figs/s_g2_2.0_Ny_5.png}
        \includegraphics[width=0.32\linewidth]{figs/s_g2_2.25_Ny_5.png}
        \includegraphics[width=0.32\linewidth]{figs/s_g2_2.5_Ny_5.png}
    \end{minipage}
\end{figure}
\begin{figure}[!hb]\ContinuedFloat
    \centering
    \begin{minipage}{\linewidth}
        \includegraphics[width=0.32\linewidth]{figs/s_g2_2.0_Ny_6.png}
        \includegraphics[width=0.32\linewidth]{figs/s_g2_2.25_Ny_6.png}
        \includegraphics[width=0.32\linewidth]{figs/s_g2_2.5_Ny_6.png}
    \end{minipage}
\end{figure}
\begin{figure}[!hb]\ContinuedFloat
    \centering
    \begin{minipage}{\linewidth}
        \includegraphics[width=0.32\linewidth]{figs/s_g2_2.0_Ny_7.png}
        \includegraphics[width=0.32\linewidth]{figs/s_g2_2.25_Ny_7.png}
        \includegraphics[width=0.32\linewidth]{figs/s_g2_2.5_Ny_7.png}
    \end{minipage}
\end{figure}
\begin{figure}[!hb]\ContinuedFloat
    \centering
    \begin{minipage}{\linewidth}
        \includegraphics[width=0.32\linewidth]{figs/s_g2_2.0_Ny_8.png}
        \includegraphics[width=0.32\linewidth]{figs/s_g2_2.25_Ny_8.png}
        \includegraphics[width=0.32\linewidth]{figs/s_g2_2.5_Ny_8.png}
    \end{minipage}
\end{figure}
\begin{figure}[!hb]\ContinuedFloat
    \centering
    \begin{minipage}{\linewidth}
        \includegraphics[width=0.32\linewidth]{figs/s_g2_2.0_Ny_9.png}
        \includegraphics[width=0.32\linewidth]{figs/s_g2_2.25_Ny_9.png}
        \includegraphics[width=0.32\linewidth]{figs/s_g2_2.5_Ny_9.png}
    \end{minipage}
\end{figure}
\begin{figure}[!hb]\ContinuedFloat
    \centering
    \begin{minipage}{\linewidth}
        \includegraphics[width=0.32\linewidth]{figs/s_g2_2.0_Ny_10.png}
        \includegraphics[width=0.32\linewidth]{figs/s_g2_2.25_Ny_10.png}
        \includegraphics[width=0.32\linewidth]{figs/s_g2_2.5_Ny_10.png}
    \end{minipage}
\end{figure}
\begin{figure}[!hb]\ContinuedFloat
    \centering
    \begin{minipage}{\linewidth}
        \includegraphics[width=0.32\linewidth]{figs/s_g2_2.0_Ny_11.png}
        \includegraphics[width=0.32\linewidth]{figs/s_g2_2.25_Ny_11.png}
        \includegraphics[width=0.32\linewidth]{figs/s_g2_2.5_Ny_11.png}
    \end{minipage}
\end{figure}
\begin{figure}[!hb]\ContinuedFloat
    \centering
    \begin{minipage}{\linewidth}
        \includegraphics[width=0.32\linewidth]{figs/s_g2_2.0_Ny_12.png}
        \includegraphics[width=0.32\linewidth]{figs/s_g2_2.25_Ny_12.png}
        \includegraphics[width=0.32\linewidth]{figs/s_g2_2.5_Ny_12.png}
    \end{minipage}
    \caption{String cross section for $g^2=2.0,\,2.25,\,2.5$ (standard Gaussian in blue, compact Gaussian in orange)}
\end{figure}

\begin{figure}[!hb]
    \centering
    \begin{minipage}{\linewidth}
        \includegraphics[width=0.32\linewidth]{figs/s_g2_2.75_Ny_3.png}
        \includegraphics[width=0.32\linewidth]{figs/s_g2_3.0_Ny_3.png}
        \includegraphics[width=0.32\linewidth]{figs/s_g2_3.25_Ny_3.png}
    \end{minipage}
\end{figure}
\begin{figure}[!hb]\ContinuedFloat
    \centering
    \begin{minipage}{\linewidth}
        \includegraphics[width=0.32\linewidth]{figs/s_g2_2.75_Ny_4.png}
        \includegraphics[width=0.32\linewidth]{figs/s_g2_3.0_Ny_4.png}
        \includegraphics[width=0.32\linewidth]{figs/s_g2_3.25_Ny_4.png}
    \end{minipage}
\end{figure}
\begin{figure}[!hb]\ContinuedFloat
    \centering
    \begin{minipage}{\linewidth}
        \includegraphics[width=0.32\linewidth]{figs/s_g2_2.75_Ny_5.png}
        \includegraphics[width=0.32\linewidth]{figs/s_g2_3.0_Ny_5.png}
        \includegraphics[width=0.32\linewidth]{figs/s_g2_3.25_Ny_5.png}
    \end{minipage}
\end{figure}
\begin{figure}[!hb]\ContinuedFloat
    \centering
    \begin{minipage}{\linewidth}
        \includegraphics[width=0.32\linewidth]{figs/s_g2_2.75_Ny_6.png}
        \includegraphics[width=0.32\linewidth]{figs/s_g2_3.0_Ny_6.png}
        \includegraphics[width=0.32\linewidth]{figs/s_g2_3.25_Ny_6.png}
    \end{minipage}
\end{figure}
\begin{figure}[!hb]\ContinuedFloat
    \centering
    \begin{minipage}{\linewidth}
        \includegraphics[width=0.32\linewidth]{figs/s_g2_2.75_Ny_7.png}
        \includegraphics[width=0.32\linewidth]{figs/s_g2_3.0_Ny_7.png}
        \includegraphics[width=0.32\linewidth]{figs/s_g2_3.25_Ny_7.png}
    \end{minipage}
\end{figure}
\begin{figure}[!hb]\ContinuedFloat
    \centering
    \begin{minipage}{\linewidth}
        \includegraphics[width=0.32\linewidth]{figs/s_g2_2.75_Ny_8.png}
        \includegraphics[width=0.32\linewidth]{figs/s_g2_3.0_Ny_8.png}
        \includegraphics[width=0.32\linewidth]{figs/s_g2_3.25_Ny_8.png}
    \end{minipage}
\end{figure}
\begin{figure}[!hb]\ContinuedFloat
    \centering
    \begin{minipage}{\linewidth}
        \includegraphics[width=0.32\linewidth]{figs/s_g2_2.75_Ny_9.png}
        \includegraphics[width=0.32\linewidth]{figs/s_g2_3.0_Ny_9.png}
        \includegraphics[width=0.32\linewidth]{figs/s_g2_3.25_Ny_9.png}
    \end{minipage}
\end{figure}
\begin{figure}[!hb]\ContinuedFloat
    \centering
    \begin{minipage}{\linewidth}
        \includegraphics[width=0.32\linewidth]{figs/s_g2_2.75_Ny_10.png}
        \includegraphics[width=0.32\linewidth]{figs/s_g2_3.0_Ny_10.png}
        \includegraphics[width=0.32\linewidth]{figs/s_g2_3.25_Ny_10.png}
    \end{minipage}
\end{figure}
\begin{figure}[!hb]\ContinuedFloat
    \centering
    \begin{minipage}{\linewidth}
        \includegraphics[width=0.32\linewidth]{figs/s_g2_2.75_Ny_11.png}
        \includegraphics[width=0.32\linewidth]{figs/s_g2_3.0_Ny_11.png}
        \includegraphics[width=0.32\linewidth]{figs/s_g2_3.25_Ny_11.png}
    \end{minipage}
\end{figure}
\begin{figure}[!hb]\ContinuedFloat
    \centering
    \begin{minipage}{\linewidth}
        \includegraphics[width=0.32\linewidth]{figs/s_g2_2.75_Ny_12.png}
        \includegraphics[width=0.32\linewidth]{figs/s_g2_3.0_Ny_12.png}
        \includegraphics[width=0.32\linewidth]{figs/s_g2_3.25_Ny_12.png}
    \end{minipage}
    \caption{String cross section for $g^2=2.75,\,3.0,\,3.25$ (standard Gaussian in blue, compact Gaussian in orange)}
\end{figure}

\begin{figure}[!hb]
    \centering
    \begin{minipage}{\linewidth}
        \includegraphics[width=0.32\linewidth]{figs/s_g2_3.5_Ny_3.png}
        \includegraphics[width=0.32\linewidth]{figs/s_g2_3.75_Ny_3.png}
        \includegraphics[width=0.32\linewidth]{figs/s_g2_4.0_Ny_3.png}
    \end{minipage}
\end{figure}
\begin{figure}[!hb]\ContinuedFloat
    \centering
    \begin{minipage}{\linewidth}
        \includegraphics[width=0.32\linewidth]{figs/s_g2_3.5_Ny_4.png}
        \includegraphics[width=0.32\linewidth]{figs/s_g2_3.75_Ny_4.png}
        \includegraphics[width=0.32\linewidth]{figs/s_g2_4.0_Ny_4.png}
    \end{minipage}
\end{figure}
\begin{figure}[!hb]\ContinuedFloat
    \centering
    \begin{minipage}{\linewidth}
        \includegraphics[width=0.32\linewidth]{figs/s_g2_3.5_Ny_5.png}
        \includegraphics[width=0.32\linewidth]{figs/s_g2_3.75_Ny_5.png}
        \includegraphics[width=0.32\linewidth]{figs/s_g2_4.0_Ny_5.png}
    \end{minipage}
\end{figure}
\begin{figure}[!hb]\ContinuedFloat
    \centering
    \begin{minipage}{\linewidth}
        \includegraphics[width=0.32\linewidth]{figs/s_g2_3.5_Ny_6.png}
        \includegraphics[width=0.32\linewidth]{figs/s_g2_3.75_Ny_6.png}
        \includegraphics[width=0.32\linewidth]{figs/s_g2_4.0_Ny_6.png}
    \end{minipage}
\end{figure}
\begin{figure}[!hb]\ContinuedFloat
    \centering
    \begin{minipage}{\linewidth}
        \includegraphics[width=0.32\linewidth]{figs/s_g2_3.5_Ny_7.png}
        \includegraphics[width=0.32\linewidth]{figs/s_g2_3.75_Ny_7.png}
        \includegraphics[width=0.32\linewidth]{figs/s_g2_4.0_Ny_7.png}
    \end{minipage}
\end{figure}
\begin{figure}[!hb]\ContinuedFloat
    \centering
    \begin{minipage}{\linewidth}
        \includegraphics[width=0.32\linewidth]{figs/s_g2_3.5_Ny_8.png}
        \includegraphics[width=0.32\linewidth]{figs/s_g2_3.75_Ny_8.png}
        \includegraphics[width=0.32\linewidth]{figs/s_g2_4.0_Ny_8.png}
    \end{minipage}
\end{figure}
\begin{figure}[!hb]\ContinuedFloat
    \centering
    \begin{minipage}{\linewidth}
        \includegraphics[width=0.32\linewidth]{figs/s_g2_3.5_Ny_9.png}
        \includegraphics[width=0.32\linewidth]{figs/s_g2_3.75_Ny_9.png}
        \includegraphics[width=0.32\linewidth]{figs/s_g2_4.0_Ny_9.png}
    \end{minipage}
\end{figure}
\begin{figure}[!hb]\ContinuedFloat
    \centering
    \begin{minipage}{\linewidth}
        \includegraphics[width=0.32\linewidth]{figs/s_g2_3.5_Ny_10.png}
        \includegraphics[width=0.32\linewidth]{figs/s_g2_3.75_Ny_10.png}
        \includegraphics[width=0.32\linewidth]{figs/s_g2_4.0_Ny_10.png}
    \end{minipage}
\end{figure}
\begin{figure}[!hb]\ContinuedFloat
    \centering
    \begin{minipage}{\linewidth}
        \includegraphics[width=0.32\linewidth]{figs/s_g2_3.5_Ny_11.png}
        \includegraphics[width=0.32\linewidth]{figs/s_g2_3.75_Ny_11.png}
        \includegraphics[width=0.32\linewidth]{figs/s_g2_4.0_Ny_11.png}
    \end{minipage}
\end{figure}
\begin{figure}[!hb]\ContinuedFloat
    \centering
    \begin{minipage}{\linewidth}
        \includegraphics[width=0.32\linewidth]{figs/s_g2_3.5_Ny_12.png}
        \includegraphics[width=0.32\linewidth]{figs/s_g2_3.75_Ny_12.png}
        \includegraphics[width=0.32\linewidth]{figs/s_g2_4.0_Ny_12.png}
    \end{minipage}
    \caption{String cross section for $g^2=3.5,\,3.75,\,4.0$ (standard Gaussian in blue, compact Gaussian in orange)}
\end{figure}

\begin{figure}[!hb]
    \centering
    \begin{minipage}{\linewidth}
        \includegraphics[width=0.32\linewidth]{figs/s_g2_4.25_Ny_3.png}
        \includegraphics[width=0.32\linewidth]{figs/s_g2_4.5_Ny_3.png}
        \includegraphics[width=0.32\linewidth]{figs/s_g2_4.75_Ny_3.png}
    \end{minipage}
\end{figure}
\begin{figure}[!hb]\ContinuedFloat
    \centering
    \begin{minipage}{\linewidth}
        \includegraphics[width=0.32\linewidth]{figs/s_g2_4.25_Ny_4.png}
        \includegraphics[width=0.32\linewidth]{figs/s_g2_4.5_Ny_4.png}
        \includegraphics[width=0.32\linewidth]{figs/s_g2_4.75_Ny_4.png}
    \end{minipage}
\end{figure}
\begin{figure}[!hb]\ContinuedFloat
    \centering
    \begin{minipage}{\linewidth}
        \includegraphics[width=0.32\linewidth]{figs/s_g2_4.25_Ny_5.png}
        \includegraphics[width=0.32\linewidth]{figs/s_g2_4.5_Ny_5.png}
        \includegraphics[width=0.32\linewidth]{figs/s_g2_4.75_Ny_5.png}
    \end{minipage}
\end{figure}
\begin{figure}[!hb]\ContinuedFloat
    \centering
    \begin{minipage}{\linewidth}
        \includegraphics[width=0.32\linewidth]{figs/s_g2_4.25_Ny_6.png}
        \includegraphics[width=0.32\linewidth]{figs/s_g2_4.5_Ny_6.png}
        \includegraphics[width=0.32\linewidth]{figs/s_g2_4.75_Ny_6.png}
    \end{minipage}
\end{figure}
\begin{figure}[!hb]\ContinuedFloat
    \centering
    \begin{minipage}{\linewidth}
        \includegraphics[width=0.32\linewidth]{figs/s_g2_4.25_Ny_7.png}
        \includegraphics[width=0.32\linewidth]{figs/s_g2_4.5_Ny_7.png}
        \includegraphics[width=0.32\linewidth]{figs/s_g2_4.75_Ny_7.png}
    \end{minipage}
\end{figure}
\begin{figure}[!hb]\ContinuedFloat
    \centering
    \begin{minipage}{\linewidth}
        \includegraphics[width=0.32\linewidth]{figs/s_g2_4.25_Ny_8.png}
        \includegraphics[width=0.32\linewidth]{figs/s_g2_4.5_Ny_8.png}
        \includegraphics[width=0.32\linewidth]{figs/s_g2_4.75_Ny_8.png}
    \end{minipage}
\end{figure}
\begin{figure}[!hb]\ContinuedFloat
    \centering
    \begin{minipage}{\linewidth}
        \includegraphics[width=0.32\linewidth]{figs/s_g2_4.25_Ny_9.png}
        \includegraphics[width=0.32\linewidth]{figs/s_g2_4.5_Ny_9.png}
        \includegraphics[width=0.32\linewidth]{figs/s_g2_4.75_Ny_9.png}
    \end{minipage}
\end{figure}
\begin{figure}[!hb]\ContinuedFloat
    \centering
    \begin{minipage}{\linewidth}
        \includegraphics[width=0.32\linewidth]{figs/s_g2_4.25_Ny_10.png}
        \includegraphics[width=0.32\linewidth]{figs/s_g2_4.5_Ny_10.png}
        \includegraphics[width=0.32\linewidth]{figs/s_g2_4.75_Ny_10.png}
    \end{minipage}
\end{figure}
\begin{figure}[!hb]\ContinuedFloat
    \centering
    \begin{minipage}{\linewidth}
        \includegraphics[width=0.32\linewidth]{figs/s_g2_4.25_Ny_11.png}
        \includegraphics[width=0.32\linewidth]{figs/s_g2_4.5_Ny_11.png}
        \includegraphics[width=0.32\linewidth]{figs/s_g2_4.75_Ny_11.png}
    \end{minipage}
\end{figure}
\begin{figure}[!hb]\ContinuedFloat
    \centering
    \begin{minipage}{\linewidth}
        \includegraphics[width=0.32\linewidth]{figs/s_g2_4.25_Ny_12.png}
        \includegraphics[width=0.32\linewidth]{figs/s_g2_4.5_Ny_12.png}
        \includegraphics[width=0.32\linewidth]{figs/s_g2_4.75_Ny_12.png}
    \end{minipage}
    \caption{String cross section for $g^2=4.25,\,4.5,\,4.75$ (standard Gaussian in blue, compact Gaussian in orange)}
\end{figure}

\begin{figure}[!hb]
    \centering
    \begin{minipage}{\linewidth}
        \includegraphics[width=0.32\linewidth]{figs/s_g2_5.0_Ny_3.png}
        \includegraphics[width=0.32\linewidth]{figs/s_g2_5.25_Ny_3.png}
        \includegraphics[width=0.32\linewidth]{figs/s_g2_5.5_Ny_3.png}
    \end{minipage}
\end{figure}
\begin{figure}[!hb]\ContinuedFloat
    \centering
    \begin{minipage}{\linewidth}
        \includegraphics[width=0.32\linewidth]{figs/s_g2_5.0_Ny_4.png}
        \includegraphics[width=0.32\linewidth]{figs/s_g2_5.25_Ny_4.png}
        \includegraphics[width=0.32\linewidth]{figs/s_g2_5.5_Ny_4.png}
    \end{minipage}
\end{figure}
\begin{figure}[!hb]\ContinuedFloat
    \centering
    \begin{minipage}{\linewidth}
        \includegraphics[width=0.32\linewidth]{figs/s_g2_5.0_Ny_5.png}
        \includegraphics[width=0.32\linewidth]{figs/s_g2_5.25_Ny_5.png}
        \includegraphics[width=0.32\linewidth]{figs/s_g2_5.5_Ny_5.png}
    \end{minipage}
\end{figure}
\begin{figure}[!hb]\ContinuedFloat
    \centering
    \begin{minipage}{\linewidth}
        \includegraphics[width=0.32\linewidth]{figs/s_g2_5.0_Ny_6.png}
        \includegraphics[width=0.32\linewidth]{figs/s_g2_5.25_Ny_6.png}
        \includegraphics[width=0.32\linewidth]{figs/s_g2_5.5_Ny_6.png}
    \end{minipage}
\end{figure}
\begin{figure}[!hb]\ContinuedFloat
    \centering
    \begin{minipage}{\linewidth}
        \includegraphics[width=0.32\linewidth]{figs/s_g2_5.0_Ny_7.png}
        \includegraphics[width=0.32\linewidth]{figs/s_g2_5.25_Ny_7.png}
        \includegraphics[width=0.32\linewidth]{figs/s_g2_5.5_Ny_7.png}
    \end{minipage}
\end{figure}
\begin{figure}[!hb]\ContinuedFloat
    \centering
    \begin{minipage}{\linewidth}
        \includegraphics[width=0.32\linewidth]{figs/s_g2_5.0_Ny_8.png}
        \includegraphics[width=0.32\linewidth]{figs/s_g2_5.25_Ny_8.png}
        \includegraphics[width=0.32\linewidth]{figs/s_g2_5.5_Ny_8.png}
    \end{minipage}
\end{figure}
\begin{figure}[!hb]\ContinuedFloat
    \centering
    \begin{minipage}{\linewidth}
        \includegraphics[width=0.32\linewidth]{figs/s_g2_5.0_Ny_9.png}
        \includegraphics[width=0.32\linewidth]{figs/s_g2_5.25_Ny_9.png}
        \includegraphics[width=0.32\linewidth]{figs/s_g2_5.5_Ny_9.png}
    \end{minipage}
\end{figure}
\begin{figure}[!hb]\ContinuedFloat
    \centering
    \begin{minipage}{\linewidth}
        \includegraphics[width=0.32\linewidth]{figs/s_g2_5.0_Ny_10.png}
        \includegraphics[width=0.32\linewidth]{figs/s_g2_5.25_Ny_10.png}
        \includegraphics[width=0.32\linewidth]{figs/s_g2_5.5_Ny_10.png}
    \end{minipage}
\end{figure}
\begin{figure}[!hb]\ContinuedFloat
    \centering
    \begin{minipage}{\linewidth}
        \includegraphics[width=0.32\linewidth]{figs/s_g2_5.0_Ny_11.png}
        \includegraphics[width=0.32\linewidth]{figs/s_g2_5.25_Ny_11.png}
        \includegraphics[width=0.32\linewidth]{figs/s_g2_5.5_Ny_11.png}
    \end{minipage}
\end{figure}
\begin{figure}[!hb]\ContinuedFloat
    \centering
    \begin{minipage}{\linewidth}
        \includegraphics[width=0.32\linewidth]{figs/s_g2_5.0_Ny_12.png}
        \includegraphics[width=0.32\linewidth]{figs/s_g2_5.25_Ny_12.png}
        \includegraphics[width=0.32\linewidth]{figs/s_g2_5.5_Ny_12.png}
    \end{minipage}
    \caption{String cross section for $g^2=5.0,\,5.25,\,5.5$ (standard Gaussian in blue, compact Gaussian in orange)}
\end{figure}

\begin{figure}[!hb]
    \centering
    \begin{minipage}{\linewidth}
        \includegraphics[width=0.32\linewidth]{figs/s_g2_5.75_Ny_3.png}
        \includegraphics[width=0.32\linewidth]{figs/s_g2_6.0_Ny_3.png}
        \includegraphics[width=0.32\linewidth]{figs/s_g2_6.25_Ny_3.png}
    \end{minipage}
\end{figure}
\begin{figure}[!hb]\ContinuedFloat
    \centering
    \begin{minipage}{\linewidth}
        \includegraphics[width=0.32\linewidth]{figs/s_g2_5.75_Ny_4.png}
        \includegraphics[width=0.32\linewidth]{figs/s_g2_6.0_Ny_4.png}
        \includegraphics[width=0.32\linewidth]{figs/s_g2_6.25_Ny_4.png}
    \end{minipage}
\end{figure}
\begin{figure}[!hb]\ContinuedFloat
    \centering
    \begin{minipage}{\linewidth}
        \includegraphics[width=0.32\linewidth]{figs/s_g2_5.75_Ny_5.png}
        \includegraphics[width=0.32\linewidth]{figs/s_g2_6.0_Ny_5.png}
        \includegraphics[width=0.32\linewidth]{figs/s_g2_6.25_Ny_5.png}
    \end{minipage}
\end{figure}
\begin{figure}[!hb]\ContinuedFloat
    \centering
    \begin{minipage}{\linewidth}
        \includegraphics[width=0.32\linewidth]{figs/s_g2_5.75_Ny_6.png}
        \includegraphics[width=0.32\linewidth]{figs/s_g2_6.0_Ny_6.png}
        \includegraphics[width=0.32\linewidth]{figs/s_g2_6.25_Ny_6.png}
    \end{minipage}
\end{figure}
\begin{figure}[!hb]\ContinuedFloat
    \centering
    \begin{minipage}{\linewidth}
        \includegraphics[width=0.32\linewidth]{figs/s_g2_5.75_Ny_7.png}
        \includegraphics[width=0.32\linewidth]{figs/s_g2_6.0_Ny_7.png}
        \includegraphics[width=0.32\linewidth]{figs/s_g2_6.25_Ny_7.png}
    \end{minipage}
\end{figure}
\begin{figure}[!hb]\ContinuedFloat
    \centering
    \begin{minipage}{\linewidth}
        \includegraphics[width=0.32\linewidth]{figs/s_g2_5.75_Ny_8.png}
        \includegraphics[width=0.32\linewidth]{figs/s_g2_6.0_Ny_8.png}
        \includegraphics[width=0.32\linewidth]{figs/s_g2_6.25_Ny_8.png}
    \end{minipage}
\end{figure}
\begin{figure}[!hb]\ContinuedFloat
    \centering
    \begin{minipage}{\linewidth}
        \includegraphics[width=0.32\linewidth]{figs/s_g2_5.75_Ny_9.png}
        \includegraphics[width=0.32\linewidth]{figs/s_g2_6.0_Ny_9.png}
        \includegraphics[width=0.32\linewidth]{figs/s_g2_6.25_Ny_9.png}
    \end{minipage}
\end{figure}
\begin{figure}[!hb]\ContinuedFloat
    \centering
    \begin{minipage}{\linewidth}
        \includegraphics[width=0.32\linewidth]{figs/s_g2_5.75_Ny_10.png}
        \includegraphics[width=0.32\linewidth]{figs/s_g2_6.0_Ny_10.png}
        \includegraphics[width=0.32\linewidth]{figs/s_g2_6.25_Ny_10.png}
    \end{minipage}
\end{figure}
\begin{figure}[!hb]\ContinuedFloat
    \centering
    \begin{minipage}{\linewidth}
        \includegraphics[width=0.32\linewidth]{figs/s_g2_5.75_Ny_11.png}
        \includegraphics[width=0.32\linewidth]{figs/s_g2_6.0_Ny_11.png}
        \includegraphics[width=0.32\linewidth]{figs/s_g2_6.25_Ny_11.png}
    \end{minipage}
\end{figure}
\begin{figure}[!hb]\ContinuedFloat
    \centering
    \begin{minipage}{\linewidth}
        \includegraphics[width=0.32\linewidth]{figs/s_g2_5.75_Ny_12.png}
        \includegraphics[width=0.32\linewidth]{figs/s_g2_6.0_Ny_12.png}
        \includegraphics[width=0.32\linewidth]{figs/s_g2_6.25_Ny_12.png}
    \end{minipage}
    \caption{String cross section for $g^2=5.75,\,6.0,\,6.25$ (standard Gaussian in blue, compact Gaussian in orange)}
\end{figure}

\begin{figure}[!hb]
    \centering
    \begin{minipage}{\linewidth}
        \includegraphics[width=0.32\linewidth]{figs/s_g2_6.5_Ny_3.png}
        \includegraphics[width=0.32\linewidth]{figs/s_g2_6.75_Ny_3.png}
        \includegraphics[width=0.32\linewidth]{figs/s_g2_7.0_Ny_3.png}
    \end{minipage}
\end{figure}
\begin{figure}[!hb]\ContinuedFloat
    \centering
    \begin{minipage}{\linewidth}
        \includegraphics[width=0.32\linewidth]{figs/s_g2_6.5_Ny_4.png}
        \includegraphics[width=0.32\linewidth]{figs/s_g2_6.75_Ny_4.png}
        \includegraphics[width=0.32\linewidth]{figs/s_g2_7.0_Ny_4.png}
    \end{minipage}
\end{figure}
\begin{figure}[!hb]\ContinuedFloat
    \centering
    \begin{minipage}{\linewidth}
        \includegraphics[width=0.32\linewidth]{figs/s_g2_6.5_Ny_5.png}
        \includegraphics[width=0.32\linewidth]{figs/s_g2_6.75_Ny_5.png}
        \includegraphics[width=0.32\linewidth]{figs/s_g2_7.0_Ny_5.png}
    \end{minipage}
\end{figure}
\begin{figure}[!hb]\ContinuedFloat
    \centering
    \begin{minipage}{\linewidth}
        \includegraphics[width=0.32\linewidth]{figs/s_g2_6.5_Ny_6.png}
        \includegraphics[width=0.32\linewidth]{figs/s_g2_6.75_Ny_6.png}
        \includegraphics[width=0.32\linewidth]{figs/s_g2_7.0_Ny_6.png}
    \end{minipage}
\end{figure}
\begin{figure}[!hb]\ContinuedFloat
    \centering
    \begin{minipage}{\linewidth}
        \includegraphics[width=0.32\linewidth]{figs/s_g2_6.5_Ny_7.png}
        \includegraphics[width=0.32\linewidth]{figs/s_g2_6.75_Ny_7.png}
        \includegraphics[width=0.32\linewidth]{figs/s_g2_7.0_Ny_7.png}
    \end{minipage}
\end{figure}
\begin{figure}[!hb]\ContinuedFloat
    \centering
    \begin{minipage}{\linewidth}
        \includegraphics[width=0.32\linewidth]{figs/s_g2_6.5_Ny_8.png}
        \includegraphics[width=0.32\linewidth]{figs/s_g2_6.75_Ny_8.png}
        \includegraphics[width=0.32\linewidth]{figs/s_g2_7.0_Ny_8.png}
    \end{minipage}
\end{figure}
\begin{figure}[!hb]\ContinuedFloat
    \centering
    \begin{minipage}{\linewidth}
        \includegraphics[width=0.32\linewidth]{figs/s_g2_6.5_Ny_9.png}
        \includegraphics[width=0.32\linewidth]{figs/s_g2_6.75_Ny_9.png}
        \includegraphics[width=0.32\linewidth]{figs/s_g2_7.0_Ny_9.png}
    \end{minipage}
\end{figure}
\begin{figure}[!hb]\ContinuedFloat
    \centering
    \begin{minipage}{\linewidth}
        \includegraphics[width=0.32\linewidth]{figs/s_g2_6.5_Ny_10.png}
        \includegraphics[width=0.32\linewidth]{figs/s_g2_6.75_Ny_10.png}
        \includegraphics[width=0.32\linewidth]{figs/s_g2_7.0_Ny_10.png}
    \end{minipage}
\end{figure}
\begin{figure}[!hb]\ContinuedFloat
    \centering
    \begin{minipage}{\linewidth}
        \includegraphics[width=0.32\linewidth]{figs/s_g2_6.5_Ny_11.png}
        \includegraphics[width=0.32\linewidth]{figs/s_g2_6.75_Ny_11.png}
        \includegraphics[width=0.32\linewidth]{figs/s_g2_7.0_Ny_11.png}
    \end{minipage}
\end{figure}
\begin{figure}[!hb]\ContinuedFloat
    \centering
    \begin{minipage}{\linewidth}
        \includegraphics[width=0.32\linewidth]{figs/s_g2_6.5_Ny_12.png}
        \includegraphics[width=0.32\linewidth]{figs/s_g2_6.75_Ny_12.png}
        \includegraphics[width=0.32\linewidth]{figs/s_g2_7.0_Ny_12.png}
    \end{minipage}
    \caption{String cross section for $g^2=6.5,\,6.75,\,7.0$ (standard Gaussian in blue, compact Gaussian in orange)}
\end{figure}

\begin{figure}[!hb]
    \centering
    \begin{minipage}{\linewidth}
        \includegraphics[width=0.32\linewidth]{figs/s_g2_7.25_Ny_3.png}
        \includegraphics[width=0.32\linewidth]{figs/s_g2_7.5_Ny_3.png}
        \includegraphics[width=0.32\linewidth]{figs/s_g2_7.75_Ny_3.png}
    \end{minipage}
\end{figure}
\begin{figure}[!hb]\ContinuedFloat
    \centering
    \begin{minipage}{\linewidth}
        \includegraphics[width=0.32\linewidth]{figs/s_g2_7.25_Ny_4.png}
        \includegraphics[width=0.32\linewidth]{figs/s_g2_7.5_Ny_4.png}
        \includegraphics[width=0.32\linewidth]{figs/s_g2_7.75_Ny_4.png}
    \end{minipage}
\end{figure}
\begin{figure}[!hb]\ContinuedFloat
    \centering
    \begin{minipage}{\linewidth}
        \includegraphics[width=0.32\linewidth]{figs/s_g2_7.25_Ny_5.png}
        \includegraphics[width=0.32\linewidth]{figs/s_g2_7.5_Ny_5.png}
        \includegraphics[width=0.32\linewidth]{figs/s_g2_7.75_Ny_5.png}
    \end{minipage}
\end{figure}
\begin{figure}[!hb]\ContinuedFloat
    \centering
    \begin{minipage}{\linewidth}
        \includegraphics[width=0.32\linewidth]{figs/s_g2_7.25_Ny_6.png}
        \includegraphics[width=0.32\linewidth]{figs/s_g2_7.5_Ny_6.png}
        \includegraphics[width=0.32\linewidth]{figs/s_g2_7.75_Ny_6.png}
    \end{minipage}
\end{figure}
\begin{figure}[!hb]\ContinuedFloat
    \centering
    \begin{minipage}{\linewidth}
        \includegraphics[width=0.32\linewidth]{figs/s_g2_7.25_Ny_7.png}
        \includegraphics[width=0.32\linewidth]{figs/s_g2_7.5_Ny_7.png}
        \includegraphics[width=0.32\linewidth]{figs/s_g2_7.75_Ny_7.png}
    \end{minipage}
\end{figure}
\begin{figure}[!hb]\ContinuedFloat
    \centering
    \begin{minipage}{\linewidth}
        \includegraphics[width=0.32\linewidth]{figs/s_g2_7.25_Ny_8.png}
        \includegraphics[width=0.32\linewidth]{figs/s_g2_7.5_Ny_8.png}
        \includegraphics[width=0.32\linewidth]{figs/s_g2_7.75_Ny_8.png}
    \end{minipage}
\end{figure}
\begin{figure}[!hb]\ContinuedFloat
    \centering
    \begin{minipage}{\linewidth}
        \includegraphics[width=0.32\linewidth]{figs/s_g2_7.25_Ny_9.png}
        \includegraphics[width=0.32\linewidth]{figs/s_g2_7.5_Ny_9.png}
        \includegraphics[width=0.32\linewidth]{figs/s_g2_7.75_Ny_9.png}
    \end{minipage}
\end{figure}
\begin{figure}[!hb]\ContinuedFloat
    \centering
    \begin{minipage}{\linewidth}
        \includegraphics[width=0.32\linewidth]{figs/s_g2_7.25_Ny_10.png}
        \includegraphics[width=0.32\linewidth]{figs/s_g2_7.5_Ny_10.png}
        \includegraphics[width=0.32\linewidth]{figs/s_g2_7.75_Ny_10.png}
    \end{minipage}
\end{figure}
\begin{figure}[!hb]\ContinuedFloat
    \centering
    \begin{minipage}{\linewidth}
        \includegraphics[width=0.32\linewidth]{figs/s_g2_7.25_Ny_11.png}
        \includegraphics[width=0.32\linewidth]{figs/s_g2_7.5_Ny_11.png}
        \includegraphics[width=0.32\linewidth]{figs/s_g2_7.75_Ny_11.png}
    \end{minipage}
\end{figure}
\begin{figure}[!hb]\ContinuedFloat
    \centering
    \begin{minipage}{\linewidth}
        \includegraphics[width=0.32\linewidth]{figs/s_g2_7.25_Ny_12.png}
        \includegraphics[width=0.32\linewidth]{figs/s_g2_7.5_Ny_12.png}
        \includegraphics[width=0.32\linewidth]{figs/s_g2_7.75_Ny_12.png}
    \end{minipage}
    \caption{String cross section for $g^2=7.25,\,7.5,\,7.75$ (standard Gaussian in blue, compact Gaussian in orange)}
\end{figure}

\begin{figure}[!hb]
    \centering
    \includegraphics[width=0.32\linewidth]{figs/s_g2_8.0_Ny_3.png}
\end{figure}
\begin{figure}[!hb]\ContinuedFloat
    \centering
    \includegraphics[width=0.32\linewidth]{figs/s_g2_8.0_Ny_4.png}
\end{figure}
\begin{figure}[!hb]\ContinuedFloat
    \centering
    \includegraphics[width=0.32\linewidth]{figs/s_g2_8.0_Ny_5.png}
\end{figure}
\begin{figure}[!hb]\ContinuedFloat
    \centering
    \includegraphics[width=0.32\linewidth]{figs/s_g2_8.0_Ny_6.png}
\end{figure}
\begin{figure}[!hb]\ContinuedFloat
    \centering
    \includegraphics[width=0.32\linewidth]{figs/s_g2_8.0_Ny_7.png}
\end{figure}
\begin{figure}[!hb]\ContinuedFloat
    \centering
    \includegraphics[width=0.32\linewidth]{figs/s_g2_8.0_Ny_8.png}
\end{figure}
\begin{figure}[!hb]\ContinuedFloat
    \centering
    \includegraphics[width=0.32\linewidth]{figs/s_g2_8.0_Ny_9.png}
\end{figure}
\begin{figure}[!hb]\ContinuedFloat
    \centering
    \includegraphics[width=0.32\linewidth]{figs/s_g2_8.0_Ny_10.png}
\end{figure}
\begin{figure}[!hb]\ContinuedFloat
    \centering
    \includegraphics[width=0.32\linewidth]{figs/s_g2_8.0_Ny_11.png}
\end{figure}
\begin{figure}[!hb]\ContinuedFloat
    \centering
    \includegraphics[width=0.32\linewidth]{figs/s_g2_8.0_Ny_12.png}
    \caption{String cross section for $g^2=8.0$ (standard Gaussian in blue, compact Gaussian in orange)}
\end{figure}
\end{document}